% Options for packages loaded elsewhere
\PassOptionsToPackage{unicode}{hyperref}
\PassOptionsToPackage{hyphens}{url}
%
\documentclass[
]{article}
\usepackage{amsmath,amssymb}
\usepackage{iftex}
\ifPDFTeX
  \usepackage[T1]{fontenc}
  \usepackage[utf8]{inputenc}
  \usepackage{textcomp} % provide euro and other symbols
\else % if luatex or xetex
  \usepackage{unicode-math} % this also loads fontspec
  \defaultfontfeatures{Scale=MatchLowercase}
  \defaultfontfeatures[\rmfamily]{Ligatures=TeX,Scale=1}
\fi
\usepackage{lmodern}
\ifPDFTeX\else
  % xetex/luatex font selection
\fi
% Use upquote if available, for straight quotes in verbatim environments
\IfFileExists{upquote.sty}{\usepackage{upquote}}{}
\IfFileExists{microtype.sty}{% use microtype if available
  \usepackage[]{microtype}
  \UseMicrotypeSet[protrusion]{basicmath} % disable protrusion for tt fonts
}{}
\makeatletter
\@ifundefined{KOMAClassName}{% if non-KOMA class
  \IfFileExists{parskip.sty}{%
    \usepackage{parskip}
  }{% else
    \setlength{\parindent}{0pt}
    \setlength{\parskip}{6pt plus 2pt minus 1pt}}
}{% if KOMA class
  \KOMAoptions{parskip=half}}
\makeatother
\usepackage{xcolor}
\usepackage[margin=1in]{geometry}
\usepackage{graphicx}
\makeatletter
\def\maxwidth{\ifdim\Gin@nat@width>\linewidth\linewidth\else\Gin@nat@width\fi}
\def\maxheight{\ifdim\Gin@nat@height>\textheight\textheight\else\Gin@nat@height\fi}
\makeatother
% Scale images if necessary, so that they will not overflow the page
% margins by default, and it is still possible to overwrite the defaults
% using explicit options in \includegraphics[width, height, ...]{}
\setkeys{Gin}{width=\maxwidth,height=\maxheight,keepaspectratio}
% Set default figure placement to htbp
\makeatletter
\def\fps@figure{htbp}
\makeatother
\setlength{\emergencystretch}{3em} % prevent overfull lines
\providecommand{\tightlist}{%
  \setlength{\itemsep}{0pt}\setlength{\parskip}{0pt}}
\setcounter{secnumdepth}{5}
\usepackage{caption}
\usepackage{multirow}
\usepackage{tabularx}
\usepackage{booktabs}
\ifLuaTeX
  \usepackage{selnolig}  % disable illegal ligatures
\fi
\usepackage{bookmark}
\IfFileExists{xurl.sty}{\usepackage{xurl}}{} % add URL line breaks if available
\urlstyle{same}
\hypersetup{
  pdftitle={Hydroclimatic variability drove human-megafauna-environment interactions during the late Pleistocene/Early Holocene in Central Chile},
  pdfkeywords={Pleistocene-Holocene transition; Ancient Tagua Tagua
Lake; Geoarchaeology; Central Chile paleoclimate; Sedimentary Facies
Analysis; Anthropogenic Ecosystems; Biogeochemical Proxies},
  hidelinks,
  pdfcreator={LaTeX via pandoc}}

\title{Hydroclimatic variability drove human-megafauna-environment
interactions during the late Pleistocene/Early Holocene in Central
Chile}
\author{true \and true \and true \and true \and true \and true \and true \and true \and true \and true \and true \and true \and true \and true}
\date{2025-05-19}

\begin{document}
\maketitle
\begin{abstract}
Major environmental changes were occurring when the first modern humans
arrived in South America during the Pleistocene-Holocene transition. How
these changes shaped human-environmental interactions across this period
remain unclear. We analyzed the stratigraphy, biogeochemistry, and
paleoclimatic models from the Ancient TaguaTagua Lake (ATTL) in central
Chile, one of the few continuous records of human and megafauna
interactions with their environment in South America, to reconstruct its
dynamics over the past 20,000 years. Results reveal that the ATTL
transitioned from a shallow, cool lake with storm-driven alluvial
deposition to a warmer, deeper, and more productive lake about 12,500
years ago, aligning with early humans' arrival. The ATTL became wetter
but experienced severe droughts between 11,000 and 8,500 years ago,
linked to shifts in Southern Westerly Winds and ENSO-like patterns.
Fluctuating conditions drove humans and fauna to seek refuge in the
basin, emphasizing local paleohydrology's role in shaping early
human-ecosystem interactions.
\end{abstract}

\section{Introduction}\label{introduction}

Modern humans settled the Americas at a time of profound ecological and
hydroclimatic transformations (1--3). These Rapid Environmental Changes
(REC)---including the extinction and redistribution of key
ecosystem-structuring species, shifts in resource availability, and
alterations in atmospheric and oceanic circulation---represented one of
the last major continental-scale socio-ecological thresholds before
modern global environmental changes (4--7). Such transformations
reconfigured ecological landscapes and could have shaped how early human
populations adapted and expanded across South America (8).

Yet, how did these REC shape early human migration, settlement, and
adaptation during the initial colonization of South America? A range of
diverse sources of information, including palaeoecological,
archaeological, and climatic data, are needed to adequately address this
question (9). The role of climate teleconnections, including large-scale
systems such as the Southeast Pacific Subtropical Anticyclone (SPSA) and
the Southern Westerly Winds (SWW), have been identified as key factors
influencing regional hydroclimates along the western slope of the
subtropical Andes during the Last Glacial Termination (Fig. 1) (10-13).
Presently, prolonged positive anomalies in the SPSA block storm fronts
from reaching the subtropical Andes (Fig. 1a), frequently resulting in
winter drought which impacts water resources, regional livestock, and
agriculture (13). In contrast, the SWW has entrained storm systems that
travel across the southern mid-latitudes, redistributing moisture and
bringing heavy rainfall that support ecosystems but also cause flooding
and economic losses (13). These climatic drivers are modulated by
phenomena such as the El Niño--Southern Oscillation (ENSO) and the
Southern Annular Mode (SAM), which influence regional and global climate
dynamics, including the frequency and intensity of the upper-level
Pacific Subtropical Jet Stream (PSJ) and atmospheric rivers (14--18). In
mid-latitudes, anomalies in these climatic drivers are marked by extreme
events, including droughts, pluvial floods, and severe storms. Early
hunter-gatherers must have been well suited and adapted to such
variability (13,19).

Despite considerable research and progress in identifying global
paleoclimatic drivers, discerning the impact of climate on human
mobility and population density in past environments remains a
significant challenge (20,21). The paucity of closely related
high-resolution paleoenvironmental and archaeological records hinders
comparisons of whether and how climate influenced human adaptation and
resource use during this early stages of exploration and colonization,
leaving uncertainties about the impact of REC (22--24). Paleoclimate
records from southern mid-latitudes indicate an intensification of the
SWW, modulated by the effects of ENSO/SAM and the Southern Ocean
(controlling CO2 fluxes) during the last deglaciation (25--28).
Alternative hypotheses suggest that the displacement of the SWW (29,30)
was driven by shifts in the Intertropical Convergence Zone (ITCZ), in
turn influenced by teleconnections with the Northern
Hemisphere---particularly those associated with North Atlantic
ice-rafting events and reductions in thermohaline circulation---as well
as by orbital changes (31--33). In contrast, the interaction of the SPSA
with climate dynamics that directly impact human activity---such as
droughts, storms and floods associated with storm tracks---remains
poorly studied (34). Past changes in the dynamics of the SPSA and SWW
provide an opportunity to study the mechanisms that cause
paleoenvironmental changes, and to assess the influence of mid-latitude
climate variability and its interaction with human activities at an
early stage (21).

The Last Glacial Termination and the onset of the Holocene marked
pronounced environmental shifts in South America, including early human
colonization and fire use, the mass extinction of large continental
mammals (\textgreater100 kg), large temperature changes (e.g.~Antarctic
Cold Reversal; ∼14.5-12.9 cal ka BP or thousands of calibrated years
before 1950), and major precipitation anomalies (e.g.~Central Andes
Pluvial Event, CAPE I and II, between 18-14 cal ka BP and 13-9 cal ka
BP, respectively) (23,35--39). Furthermore, geological and
archaeological records clearly indicate that population expansion and
megafaunal extinctions occurred relatively synchronously with climate
variations observed during Termination I (T1) in the Northern
Hemisphere, such as the Heinrich Stadial 1 (HS1, \textasciitilde{} 18-15
cal ka BP), Oldest Dryas (17.5-14.5 cal ka BP), Bølling-Allerød
(14.7-12.9 cal ka BP), and Younger Dryas Stadial
(\textasciitilde12.9-11.7 cal ka BP) (1,40--43). However, the magnitude
and distribution of these possible teleconnections at the South American
mid-latitudes remain are not well-constrained. Hence, the study of past
environmental records, especially those linked to human activity, is of
paramount importance for understanding coupled human-environment systems
and the global, regional, and local mechanisms that drove ecosystem
changes during the Last Glacial Termination.

Here, we examine hydroclimatic and ecosystem shifts at the ATTL
geoarchaeological sites in central Chile, highlighting their importance
for understanding late Pleistocene/Early Holocene human--environment
interactions. Our study features a comprehensive multi-proxy analysis of
a newly obtained 2.8-meter stratigraphic profile from the TT19-3A
sequence (Fig. 1b), offering key insights into the climatic conditions
surrounding the arrival of the first hunter-gatherers in the region.
This record integrates data from sedimentary facies, grain size
distributions (GSDs), geochemical proxies (TOC, TIC, C/N, TS, δ13Ccar,
and δ18Ocar), and biological indicators, including phytoliths, pollen,
non-pollen palynomorphs (NPPs), and diatoms. In addition, we use an
ensemble of simulations from the Paleoclimate Modeling Intercomparison
Project phase 4 (PMIP4), integrated into the Coupled Model
Intercomparison Project phase 6 (CMIP6), to analyze zonal wind and
precipitation patterns during the Last Glacial Maximum (LGM) and the
Middle Holocene, focusing on their variability under humid and dry
conditions in central Chile. This approach provides a comparative
baseline of atmospheric circulation conditions before and after the
arrival of the first hunter-gatherers in central Chile, thus enabling a
direct contrast between these large-scale patterns and our local
sedimentary record. Overall, our results highlight the interplay between
hydroclimatic variability and human activity during the Last Glacial
Termination, shedding deeper insight into how early populations
interacted with shifting environmental conditions.

\section{Results}\label{results}

\subsection{Stratigraphy and Facies: Depositional Variability and
Geochemical
Shifts}\label{stratigraphy-and-facies-depositional-variability-and-geochemical-shifts}

The TT19-3A sequence comprises 15 sedimentary facies distributed in 8
stratigraphic units, and several sedimentary clusters (Figs. 1, 2 and
supplementary material). At the base (Unit 8; 2.80 m - 2.30 m), brown to
greenish sediments include clayey silty facies (\emph{L8; 2.80 m - 2.61
m}) and sandy facies (\emph{L7; 2.61 m - 2.33 m}), associated with
clusters 4, 7, and 2 (Fig. 2). A coarsening-upward sequence (\emph{L8b}
\rightarrow L8a \rightarrow L7) with desiccation cracks lies above these
sandy layers and may indicate subaerial exposure or subaquatic
fracturing related to degassing.

Units 7 and 6 consist of greenish-brown sandy and silty sediments
(\emph{facies L6; 2.33 m - 1.84 m; and L5; 1.84 m - 1.65 m}), linked to
clusters 3 and 7 (Figs. 2 and S1). Unit 7 forms a fining-upward sequence
with fine sands containing subrounded clasts and dispersed fish
vertebrae (\emph{L6b}), overlain by sandy silts (\emph{L6a}). Unit 6 is
composed of dark grayish-green sandy silts with reworked gastropod and
bivalve shells.

A marked erosive unconformity formed at the base of Unit 5 (1.65 m -
1.50 m), which is composed of massive light grayish green silt with
millimeter to centimeter-sized carbon fragments, faunal remains and
lithic artifacts, and abundant \emph{Biomphalaria taguataguensis}
(\emph{Planorbidae}) shells (\emph{L4b facies}). Facies \emph{L4}
corresponds to cluster 1. The top of Unit 5 (\emph{L4a}) is a dark
brown, fine-grained, carbonate-rich (6\% TIC) silt layer (Figs. 2 and
S1).

Unit 4 (1.50 m - 0.80 m) contains carbonate (2-5\% TIC) and diatom-rich
coarse gray silts and fine sands, defined by four facies. Facies
\emph{L3d} (1.50 m - 1.30 m) consists of fine-grained light gray silts
with some large volcanic clasts up to 3.5 cm. Facies \emph{L3c} (1.30 m
- 1.10 m) is characterized by massive to banded gray clay silts with
some millimetric clasts, higher carbonate content (up to 4\% TIC), and
abundant shells of \emph{B. taguataguensis}, randomly dispersed but more
abundant towards the top. Facies \emph{L3b} (1.10 m - 0.90 m) comprises
massive to banded fine-to-medium gray clay silts with \emph{B.
taguataguensis} shells and millimeter- to centimeter-sized charcoal
fragments. Facies \emph{L3a} (0.90 m - 0.80 m) are massive
medium-grained gray to dark gray sandy silts with sub-rounded clasts,
higher carbonate content and shells of \emph{B. taguataguensis}.

Units 3 and 2 contain finer gray sediments with relatively high
carbonate (Unit 3, 1-5\% TIC) and organic (Unit 2, up to 2.5\% TOC)
contents. Unit 3 (0.80 m - 0.45 m) is composed of gray fine sandy silts
with silty sand lenses or pockets, and in situ shells of \emph{B.
taguataguensis} and \emph{Diplodon} sp. (\emph{facies L2}). Unit 2 (0.45
m - 0.38 m) consists of dark gray fine sandy silts with charcoal
fragments (up to 1 cm) and gastropod shells. Finally, Unit 1 (0.38 m - 0
m) consists of light brownish-gray clays with abundant diatoms, higher
organic content and presence of charcoal (\emph{facies L2}), as well as
silty sands containing rounded gravel clasts (0.5--1 cm) and larger
charcoal fragments (0.2--1.5 cm; \emph{facies L1}). Abundant clasts,
macrocharcoal, and elevated TOC and TN in Unit 1 suggest
post-depositional reworking, soil formation, and anthropogenic
influence.

Overall, sediments are mostly silt (\textasciitilde50\%) and sand
(\textasciitilde40\%), with 5--20\% clay and \textless5\% fine gravels
(Fig. 2). The TT19-3A sequence displays a multimodal grain size
distribution (GSD). Using rEMMA, we identified five end-members
{[}modes; \emph{rEM1}: 8 ɸ, \emph{rEM2}: 4.1 ɸ, \emph{rEM3}: 2.7 ɸ,
\emph{rEM4}: 1.3 ɸ, and \emph{rEM5}: 0.47 ɸ{]} that explain
\textasciitilde78\% of the total variance (Figs. S2 and S3). The
\emph{rEM2} end-member (peak at \textasciitilde60 µm) dominates the
upper half of the record (\textasciitilde40\% variance), particularly in
diatom-rich lacustrine facies (\emph{L3a, L3b, L3c}). In contrast,
\emph{rEM1} (peak \textasciitilde4 µm; \textasciitilde18\% variance)
characterizes silty clay facies such as \emph{L8} and \emph{L6}, while
\emph{rEM3} (peak \textasciitilde160 µm; \textasciitilde15\% variance)
occurs in fine-grained, diatom-rich intervals (1.30 - 0.80 m). The
\emph{rEM4} end-member, with a bimodal distribution (\textasciitilde420
µm and \textasciitilde60 µm), and \emph{rEM5}, with a dominant coarse
sand peak (\textasciitilde720 µm), each explain \textasciitilde14\% of
the variance.

Carbonate and TOC contents vary among units. Units 8 to 6 are
essentially carbonate-free, with TC \textless0.1\%. Facies \emph{L8a}
and \emph{L8b} exhibit relatively high TS (0.1-1\%) and low TC
(\textasciitilde0.1\%), alongside δ\^{}18O\_car and δ\^{}13C\_car values
near -5‰ and -10‰, respectively. Facies \emph{L7} shows a gradual
decrease in the δ\^{}18O\_car values (cluster 2, Fig. S1), reaching the
lowest values of the sequence. Unit 7 (\emph{L6 facies}) is
characterized by two sedimentary clusters: the cluster 7, with
increasing δ\^{}18O\_car and δ\^{}13C\_car values and cluster 3, with
low \emph{rEM1} and \emph{rEM2} values and the absence of \emph{rEM5}.
In Unit 6 (\emph{L5 facies}), \emph{rEM5} increases while δ\^{}18O\_car,
\emph{rEM1}, and \emph{rEM2} values decrease. In Unit 5 (\emph{L4
facies}), TC, TIC, δ\^{}18O\_car, δ\^{}13C\_car, and \emph{rEM2} values
increase, whereas \emph{rEM1}, \emph{rEM4}, and \emph{rEM5} values
diminish. \emph{L4a facies} records the highest TIC (\textgreater7\%),
δ\^{}18O\_car (\textgreater-1 ‰), δ\^{}13C\_car (\textgreater1 ‰), and
\emph{rEM2} values, along with minimal \emph{rEM1}, \emph{rEM3},
\emph{rEM4}, and \emph{rEM5}.

In Unit 4 (\emph{L3 facies}), there is a gradual increase in the TOC/TN
ratio (\textgreater5), δ\^{}18O\_car (\textgreater-1 ‰), and
δ\^{}13C\_car values, along with variable GSD percentages. In Unit 3
(\emph{L2 facies}), geochemical values show reduced variability, with
high TC, TIC, TOC/TN, and \emph{rEM2} values. Finally, in Unit 2
(\emph{facies 2}), TC values reach their highest levels, primarily due
to the increase in TOC.

\subsection{Aquatic--terrestrial ecosystem proxies and community
shifts}\label{aquaticterrestrial-ecosystem-proxies-and-community-shifts}

The diatom record in the TT19-3A sequence (Figs. 3 and S4) is dominated
by benthic taxa with high specific conductance optima---such as
\emph{Anomoeoneis sculpta} (Ehrenberg) Cleve, \emph{Nitzschia amphibia}
Grunow, \emph{Epithemia turgida} (Ehrenberg) Kützing, \emph{Cymbella
cymbiformis} Bahls, and \emph{Gomphonema affine} Kützing. This overall
dominance is interrupted by an interval of increased tychoplanktonic
small fragilarioids between 1.50 and 1.20 m (\emph{L3d}, \emph{L3c}).

At the base of the sequence (\emph{Unit 8, L8b}), diatom valves are
absent. Moving upward (\emph{L8a}, Units 7 and 6), we observe the
presence of \emph{Anomoeoneis sculpta} (most abundant), \emph{N.
amphibia}, \emph{E. turgida}, \emph{G. affine}, and \emph{Ulnaria} sp.
In Unit 5 (\emph{L4b}), \emph{A. sculpta} remains the dominant taxon
(\textasciitilde84\% abundance), and diatom valves are ubiquitous,
allowing for the estimation of relative abundances.

A notable assemblage shift occurs in \emph{L4a} (Unit 5), \emph{L3d},
and \emph{L3c} (Unit 4), where \emph{A. sculpta} nearly disappears and
diatom diversity increases substantially. Here, small fragilarioids,
\emph{N. amphibia}, \emph{E. turgida}, \emph{C. cymbiformis}, and
\emph{G. affine} all become more common, accompanied by planktonic
species such as \emph{Stephanocyclus meneghinianus} (Kütz) Kulikovskiy,
Genkal and Kociolek, \emph{Aulacoseira granulata} (Ehrenberg) Simonsen,
and \emph{Ulnaria} sp.

\emph{L4a} and \emph{L3d} differ from \emph{L3c} in that small
fragilarioids and \emph{A. granulata} are more abundant, whereas
\emph{N. amphibia}, \emph{S. meneghinianus}, and \emph{A. sculpta}
(\textasciitilde5\% abundance) are more common in \emph{L3c}.
Subsequently, \emph{L3b} records a marked recovery of \emph{A. sculpta}
(\textasciitilde77\%), accompanied by a decrease in other taxa. In
\emph{L3a}, \emph{A. sculpta} declines (\textasciitilde22\%) as \emph{G.
affine} and \emph{N. amphibia} increase.

In Unit 3 (\emph{L2}) and Unit 2 (\emph{L1}), \emph{A. sculpta}
partially recovers (\textasciitilde36\%) alongside slight increases in
planktonic taxa \emph{S. meneghinianus} and \emph{A. granulata}.

Pollen and NPPs data (Fig. 3, S5) reveal distinct assemblages of
hygro-hydrophytes, ferns, freshwater algae, and spores of coprophilous
fungi. Two well-defined sections are apparent:

\begin{enumerate}
\def\labelenumi{\arabic{enumi}.}
\tightlist
\item
  \textbf{The first section (2.63 m - 1.50 m; Units 8, 7, 6, and 5):}\\
  This section is characterized by high abundances of freshwater algae
  \emph{Botryococcus} and \emph{Spirogyra}, along with occasional
  \emph{Mougeotia} and \emph{Pediastrum}.

  \begin{itemize}
  \tightlist
  \item
    In Units 8 and 7 (2.63 m - 1.85 m), \emph{Typha}, ferns
    (\emph{Pteris}), \emph{Botryococcus}, and \emph{Spirogyra}
    dominate.\\
  \item
    Towards the upper part of this section (2.38 m - 1.50 m), the
    abundance of \emph{Botryococcus}, \emph{Spirogyra}, and arbuscular
    mycorrhizal fungi (\emph{Glomus}) remains elevated, but littoral and
    floating vegetation reach their lowest values.
  \end{itemize}
\item
  \textbf{The second section (1.50 m - 0.38 m):}\\
  This section shows a progressive increase in water surface indicators
  and the colonization of littoral (\emph{Cyperaceae}, \emph{Juncaceae},
  \emph{Typha}) and floating vegetation (\emph{Lemna}, \emph{Nuphar}),
  especially pronounced between 0.55 m and 0.38 m. Freshwater algae
  remain present throughout this interval.
\end{enumerate}

Phytolith analysis (Figs. 3 and S6) identifies opal phytoliths from
\emph{Cyperaceae}, \emph{Poaceae}, and \emph{Arecaceae}. Within the
\emph{Poaceae}, the subfamilies \emph{Bambusoideae},
\emph{Chloridoideae}, and \emph{Pooideae} were distinguished. We also
identified \emph{Jubaea chilensis} or Chilean palm (\emph{Arecaceae}),
as well as phytoliths from woody dicotyledons.

Unit 8 exhibits very low phytolith concentrations compared to the rest
of the record. Here, \emph{Poaceae} phytoliths dominate
(\textasciitilde90\%), with fewer \emph{Bambusoideae} and
\emph{Chloridoideae} phytoliths (\textasciitilde5\%). Unit 7 (2.30 m -
1.75 m), \emph{Bambusoideae} and \emph{Chloridoideae} reach their
highest percentages (46\% and 26\%, respectively), while
\emph{Cyperaceae} increases modestly (\textasciitilde13\%). Unit 6 (1.75
m - 1.65 m) shows the highest overall \emph{Poaceae} content, though
\emph{Bambusoideae}, \emph{Pooideae}, and \emph{Chloridoid} percentages
decrease.

From Unit 5 to Unit 3 (1.65 m - 0.40 m), the total phytolith content
rises, accompanied by the first appearance of woody and shrub
morphotypes indicative of sclerophyll forests (woody dicotyledons and
\emph{J. chilensis}). During this interval, \emph{Bambusoideae} and
\emph{Chloridoid} proportions slightly decline in Unit 3. Additionally,
\emph{Cyperaceae} levels increase again (to \textasciitilde13\%) between
Units 5 and 4 (1.50-1.42 m) before stabilizing at lower values
(\textasciitilde2\%) thereafter.

\subsection{Cultural Signals, Megafauna, and Chronological
Context}\label{cultural-signals-megafauna-and-chronological-context}

Archaeological evidence from the ATTL indicates that human occupation at
sites Taguatagua-1 (\emph{TT-1}), Taguatagua-2 (\emph{TT-2}), and
Taguatagua-3 (\emph{TT-3}) began around 12.6 - 11.6 ka cal BP during
deposition of Unit 5 in the TT19-3A sequence (Fig. 1c) (10, 44, 45).
These archaeological sites document direct human interactions with
megafauna---such as gomphotheres, horses, and cervids---as well as the
exploitation of smaller prey (birds, frogs) and the use of native
plants.

No archaeological radiocarbon dates have been identified between 11.6
and 9.0 ka cal BP at the ATTL; however, open-air residential camps and
extensive lithic knapping areas appear at ca. 9.0 ka cal BP. Human
burials within residential zones, forming mounds composed of dense
domestic refuse and mixed stratigraphic deposits, occur from ca. 8.2 --
5.8 ka cal BP. At least seven such mounds were identified along the
lakeshore, though only a few have been systematically excavated.
Isolated radiocarbon dates after this period suggest reduced human
presence until the onset of the Late Holocene (ca. 4.0 ka cal BP).

In a broader regional context, the summed probability distributions
(SPD) of archaeological radiocarbon dates for South America and central
Chile (32°S -- 37°S) (46), reveal a pronounced increase around 12.6 ka
cal BP. This uptick aligns with the timing of initial human presence
documented at the ATTL and the environmental shifts described earlier,
underscoring a notable rise in human activity as climates warmed and
landscapes transformed.

Throughout the Early to Mid Holocene (ca. 11.7 - 6 ka cal BP), the SPD
remains generally elevated, albeit with periods of fluctuation. These
patterns, when considered alongside the localized archaeological
evidence, suggest sustained human presence and periodic fluctuations in
occupation intensity in central Chile. The chronological framework for
the TT19-3A sequence provided by the age-depth models (\emph{Bacon} and
\emph{Bchron}; Fig. 4) supports these interpretations, offering a robust
temporal backdrop for correlating cultural and environmental shifts. The
models excluded two outlier ages. Bulk sediment and charcoal samples
taken at the same depths showed no significant differences, ruling out a
reservoir effect. According to the ensemble age-depth model, the TT19-3A
sequence covers the last 20.5 cal ka BP (Table S1 and Fig. 4), with a
higher model robustness in the upper four units. \textbf{Unit 8 (2.66 m
-- 2.30 m):} spans from \textasciitilde20.5 to 16 cal ka BP with an
accumulation rate of \textasciitilde200 a cm⁻¹. \textbf{Unit 7 (2.30 m -
1.84 m):} was deposited between \textasciitilde16 and 12.5 cal ka BP
with an accumulation rate of \textasciitilde100 a cm⁻¹. \textbf{Unit 6
(1.84 m - 1.65 m):} was deposited between 13 and 12.6 cal ka BP with an
accumulation rate of \textasciitilde20 a cm⁻¹. \textbf{Unit 5 (1.65 m -
1.50 m):} encompassing the Pleistocene-Holocene transition (12.6 - 11.3
cal ka BP), has an accumulation rate of \textasciitilde100 a cm⁻¹.
\textbf{Unit 4 (1.50 m - 0.80 m):} extends from 11.3 to 8.6 cal ka BP,
characterized by \textasciitilde50 a cm⁻¹. \textbf{Unit 3 (0.80 m - 0.45
m):} spans between 8.6 and 6.5 cal ka BP with a median rate of
\textasciitilde60 a cm⁻¹. \textbf{Unit 2 (0.45 - 0.38 m):} covers
between 6.5 and 6.0 cal ka BP with an accumulation rate of
\textasciitilde80 a cm⁻¹.

Taken together, the localized archaeological evidence from the ATTL, the
broader SPD patterns for central Chile and South America, and the robust
ensemble age-depth models yield a coherent narrative of increasing human
presence and activity beginning around the Pleistocene-Holocene
transition.

\begin{verbatim}
## > reading /home/mat150604/MEGA/MEGAsync/workflow/TT3/_targets/meta/meta
\end{verbatim}

\begin{center}\includegraphics[width=1\linewidth]{/home/mat150604/MEGA/MEGAsync/workflow/TT3/paper/paper_files/figure-latex/unnamed-chunk-3-1} \end{center}

\begin{verbatim}
## > reading /home/mat150604/MEGA/MEGAsync/workflow/TT3/_targets/meta/meta
\end{verbatim}

\begin{table}[!h]
\caption{Radiocarbon dating of TT3 sequence. The table shows the ages in terms of years `before present' (BP), percent modern carbon (pMC) and their uncertainties (1 $1\sigma$).
      The median age calibration was conducted with Oxcal 4.4 and SHCal20 calibration curve (Hogg et al., 2020; Bronk Ramsey 2009a)}
\centering
\fontsize{7}{9}\selectfont
\begin{tabular}[t]{lcccccccr}
\toprule
Lab code & Depth & $^{14}$C a BP & $1\sigma$ & pMC & $1\sigma$ & Median a cal BP & $1\sigma$ & Material\\
\midrule
D-AMS 039542 & 39 & 5345 & 28 & 51.41 & 0.18 & 6086 & 69 & Charcoal\\
D-AMS 039543 & 44 & 5806 & 33 & 48.54 & 0.20 & 6571 & 56 & Charcoal\\
D-AMS 039544 & 67 & 7259 & 32 & 40.51 & 0.16 & 8021 & 62 & Charcoal\\
D-AMS 039545 & 106 & 8657 & 45 & 34.04 & 0.19 & 9591 & 57 & Charcoal\\
D-AMS 039546 & 150 & 9736 & 35 & 29.76 & 0.13 & 11134 & 111 & Bulk sediment\\
D-AMS 039547 & 161 & 10211 & 63 & 28.05 & 0.22 & 11817 & 144 & Bulk sediment\\
D-AMS 039548 & 162 & 10936 & 41 & 25.63 & 0.13 & 12813 & 46 & Charcoal\\
D-AMS 039549 & 162 & 10747 & 49 & 26.24 & 0.16 & 12709 & 38 & Bulk sediment\\
D-AMS 039550 & 168 & 10590 & 39 & 26.76 & 0.13 & 12549 & 64 & Charcoal\\
D-AMS 039551 & 199 & 11117 & 51 & 25.06 & 0.16 & 13015 & 61 & Bulk sediment\\
D-AMS 039552 & 223 & 9977 & 45 & 28.88 & 0.16 & 11377 & 122 & Bulk sediment\\
D-AMS 039553 & 254 & 17353 & 84 & 11.53 & 0.12 & 20890 & 124 & Bulk sediment\\
D-AMS 039554 & 266 & 16558 & 69 & 12.73 & 0.11 & 19959 & 130 & Bulk sediment\\
\bottomrule
\end{tabular}
\end{table}

\begin{verbatim}
## > reading /home/mat150604/MEGA/MEGAsync/workflow/TT3/_targets/meta/meta
\end{verbatim}

\begin{verbatim}
## $core
## [1] "TT3"
## 
## $thick
## [1] 2
## 
## $dets
##          labID   age error depth cc delta.R delta.STD t.a t.b
## 1  D-AMS039542  5345    28    39  3       0         0   3   4
## 2  D-AMS039543  5806    33    44  3       0         0   3   4
## 3  D-AMS039544  7259    32    67  3       0         0   3   4
## 4  D-AMS039545  8657    45   106  3       0         0   3   4
## 5  D-AMS039546  9736    35   150  3       0         0   3   4
## 6  D-AMS039547 10211    63   161  3       0         0   3   4
## 7  D-AMS039548 10936    41   162  3       0         0   3   4
## 8  D-AMS039549 10747    49   162  3       0         0   3   4
## 9  D-AMS039550 10590    39   168  3       0         0   3   4
## 10 D-AMS039551 11117    51   199  3       0         0   3   4
## 11 D-AMS039552  9977    45   223  3       0         0   3   4
## 12 D-AMS039553 17353    84   254  3       0         0   3   4
## 13 D-AMS039554 16558    69   266  3       0         0   3   4
## 
## $d.min
## [1] 39
## 
## $d.max
## [1] 266
## 
## $d.by
## [1] 1
## 
## $depths.file
## [1] 0
## 
## $slump
##      [,1] [,2]
## [1,]  176  180
## 
## $acc.mean
## [1] 50 80
## 
## $acc.shape
## [1] 1.4 1.0
## 
## $mem.mean
## [1] 0.5
## 
## $mem.strength
## [1] 10
## 
## $boundary
## [1] NA
## 
## $hiatus.depths
## [1] 176
## 
## $hiatus.max
## [1] 10000
## 
## $BCAD
## [1] 0
## 
## $cc
## [1] 1
## 
## $postbomb
## [1] 0
## 
## $cc1
## [1] "IntCal20"
## 
## $cc2
## [1] "Marine20"
## 
## $cc3
## [1] "SHCal20"
## 
## $cc4
## [1] "ConstCal"
## 
## $depth.unit
## [1] cm
## 
## $unit
## [1] "cm"
## 
## $age.unit
## [1] yr
## 
## $normal
## [1] 0
## 
## $t.a
## [1] 3
## 
## $t.b
## [1] 4
## 
## $delta.R
## [1] 0
## 
## $delta.STD
## [1] 0
## 
## $prob
## [1] 0.95
## 
## $date
## [1] "Mon May 19 18:33:55 2025"
## 
## $runname
## [1] ""
## 
## $ssize
## [1] 8000
## 
## $dark
## [1] 1
## 
## $youngest.age
## [1] -75
## 
## $oldest.age
## [1] 1e+06
## 
## $cutoff
## [1] 0.01
## 
## $age.res
## [1] 200
## 
## $after
## [1] 5e-05
## 
## $save.info
## [1] TRUE
## 
## $coredir
## [1] "/home/mat150604/MEGA/MEGAsync/workflow/TT3/cores_bacon/"
## 
## $seed
## [1] 128442
## 
## $isplum
## [1] FALSE
## 
## $calib
## $calib$d
##  [1]  39  44  67 106 150 161 162 162 168 199 223 254 266
## 
## $calib$cc
##  [1] 3 3 3 3 3 3 3 3 3 3 3 3 3
## 
## $calib$probs
## $calib$probs[[1]]
##       [,1]         [,2]
##  [1,] 5905 0.0002733104
##  [2,] 5910 0.0004490923
##  [3,] 5915 0.0006697521
##  [4,] 5920 0.0008817601
##  [5,] 5925 0.0012801536
##  [6,] 5930 0.0019001787
##  [7,] 5935 0.0028838347
##  [8,] 5940 0.0046375168
##  [9,] 5945 0.0069111816
## [10,] 5950 0.0090799558
## [11,] 5955 0.0097906884
## [12,] 5960 0.0083540770
## [13,] 5965 0.0063572225
## [14,] 5970 0.0046375168
## [15,] 5975 0.0033497731
## [16,] 5980 0.0025613511
## [17,] 5985 0.0030773444
## [18,] 5990 0.0060088445
## [19,] 5995 0.0105173549
## [20,] 6000 0.0186437515
## [21,] 6005 0.0232429150
## [22,] 6010 0.0240712125
## [23,] 6015 0.0239575909
## [24,] 6020 0.0235385128
## [25,] 6025 0.0229842442
## [26,] 6030 0.0224811190
## [27,] 6035 0.0220334211
## [28,] 6040 0.0224255596
## [29,] 6045 0.0236965830
## [30,] 6050 0.0240679128
## [31,] 6055 0.0239371019
## [32,] 6060 0.0240864478
## [33,] 6065 0.0235927140
## [34,] 6070 0.0233506152
## [35,] 6075 0.0239314685
## [36,] 6080 0.0238689467
## [37,] 6085 0.0219631495
## [38,] 6090 0.0192128462
## [39,] 6095 0.0170714872
## [40,] 6100 0.0169787175
## [41,] 6105 0.0196438803
## [42,] 6110 0.0231330192
## [43,] 6115 0.0237102963
## [44,] 6120 0.0198991724
## [45,] 6125 0.0155891790
## [46,] 6130 0.0137895534
## [47,] 6135 0.0149966550
## [48,] 6140 0.0189301387
## [49,] 6145 0.0232128539
## [50,] 6150 0.0239371019
## [51,] 6155 0.0223679897
## [52,] 6160 0.0215243076
## [53,] 6165 0.0225567481
## [54,] 6170 0.0238689467
## [55,] 6175 0.0237887210
## [56,] 6180 0.0213836808
## [57,] 6185 0.0176581931
## [58,] 6190 0.0132783582
## [59,] 6195 0.0097589277
## [60,] 6200 0.0069111816
## [61,] 6205 0.0048570141
## [62,] 6210 0.0032350386
## [63,] 6215 0.0021339665
## [64,] 6220 0.0015789819
## [65,] 6225 0.0014003773
## [66,] 6230 0.0014869325
## [67,] 6235 0.0018845292
## [68,] 6240 0.0030621806
## [69,] 6245 0.0047747542
## [70,] 6250 0.0071710321
## [71,] 6255 0.0086779344
## [72,] 6260 0.0079776365
## [73,] 6265 0.0060851010
## [74,] 6270 0.0039696132
## [75,] 6275 0.0024821251
## [76,] 6280 0.0014290703
## [77,] 6285 0.0007460396
## [78,] 6290 0.0004372640
## [79,] 6295 0.0002616662
## 
## $calib$probs[[2]]
##       [,1]         [,2]
##  [1,] 6315 0.0003795822
##  [2,] 6320 0.0005381358
##  [3,] 6325 0.0005955125
##  [4,] 6330 0.0005519012
##  [5,] 6335 0.0005230453
##  [6,] 6340 0.0004746277
##  [7,] 6345 0.0005116987
##  [8,] 6350 0.0006387642
##  [9,] 6355 0.0006648483
## [10,] 6360 0.0007897573
## [11,] 6365 0.0007625744
## [12,] 6370 0.0005116987
## [13,] 6375 0.0003616014
## [14,] 6380 0.0002838817
## [15,] 6385 0.0002894656
## [16,] 6390 0.0004195520
## [17,] 6395 0.0007575894
## [18,] 6400 0.0013530061
## [19,] 6405 0.0023511966
## [20,] 6410 0.0039755820
## [21,] 6415 0.0048112893
## [22,] 6420 0.0041557885
## [23,] 6425 0.0031520075
## [24,] 6430 0.0026771506
## [25,] 6435 0.0027557459
## [26,] 6440 0.0039038341
## [27,] 6445 0.0057888245
## [28,] 6450 0.0074030908
## [29,] 6455 0.0088113474
## [30,] 6460 0.0090285794
## [31,] 6465 0.0088113474
## [32,] 6470 0.0084880135
## [33,] 6475 0.0076482964
## [34,] 6480 0.0076482964
## [35,] 6485 0.0093297114
## [36,] 6490 0.0130208582
## [37,] 6495 0.0186381808
## [38,] 6500 0.0234987596
## [39,] 6505 0.0252421997
## [40,] 6510 0.0254896953
## [41,] 6515 0.0246781702
## [42,] 6520 0.0233945613
## [43,] 6525 0.0227182626
## [44,] 6530 0.0228669561
## [45,] 6535 0.0241408833
## [46,] 6540 0.0248250067
## [47,] 6545 0.0248250067
## [48,] 6550 0.0251000099
## [49,] 6555 0.0263088718
## [50,] 6560 0.0272184775
## [51,] 6565 0.0265490986
## [52,] 6570 0.0250377012
## [53,] 6575 0.0240531656
## [54,] 6580 0.0251839287
## [55,] 6585 0.0266502097
## [56,] 6590 0.0272177849
## [57,] 6595 0.0272170725
## [58,] 6600 0.0272170725
## [59,] 6605 0.0268721978
## [60,] 6610 0.0265490986
## [61,] 6615 0.0266675510
## [62,] 6620 0.0266844237
## [63,] 6625 0.0258035492
## [64,] 6630 0.0238046403
## [65,] 6635 0.0205450959
## [66,] 6640 0.0180162616
## [67,] 6645 0.0175043631
## [68,] 6650 0.0163458227
## [69,] 6655 0.0150630648
## [70,] 6660 0.0124834621
## [71,] 6665 0.0090695662
## [72,] 6670 0.0059399195
## [73,] 6675 0.0036303755
## [74,] 6680 0.0026904792
## [75,] 6685 0.0023468795
## [76,] 6690 0.0025420243
## [77,] 6695 0.0029763400
## [78,] 6700 0.0033967056
## [79,] 6705 0.0037738437
## [80,] 6710 0.0040090904
## [81,] 6715 0.0043015780
## [82,] 6720 0.0043015780
## [83,] 6725 0.0038742153
## [84,] 6730 0.0030685260
## [85,] 6735 0.0024724855
## [86,] 6740 0.0016621964
## [87,] 6745 0.0009986609
## [88,] 6750 0.0006150793
## [89,] 6755 0.0004659244
## [90,] 6760 0.0005166365
## [91,] 6765 0.0006367454
## [92,] 6770 0.0007871655
## [93,] 6775 0.0007871655
## [94,] 6780 0.0006098882
## [95,] 6785 0.0005034247
## [96,] 6790 0.0003611778
## 
## $calib$probs[[3]]
##       [,1]         [,2]
##  [1,] 7845 0.0005273065
##  [2,] 7850 0.0005964623
##  [3,] 7855 0.0005656011
##  [4,] 7860 0.0007112891
##  [5,] 7865 0.0009132702
##  [6,] 7870 0.0013521219
##  [7,] 7875 0.0022666274
##  [8,] 7880 0.0030412101
##  [9,] 7885 0.0031239386
## [10,] 7890 0.0022449280
## [11,] 7895 0.0012277596
## [12,] 7900 0.0007103292
## [13,] 7905 0.0005021287
## [14,] 7910 0.0004414536
## [15,] 7915 0.0005118891
## [16,] 7920 0.0007103292
## [17,] 7925 0.0011192186
## [18,] 7930 0.0020383501
## [19,] 7935 0.0038688009
## [20,] 7940 0.0070427785
## [21,] 7945 0.0091980481
## [22,] 7950 0.0093286768
## [23,] 7955 0.0097065056
## [24,] 7960 0.0112912639
## [25,] 7965 0.0150275241
## [26,] 7970 0.0211712180
## [27,] 7975 0.0281236514
## [28,] 7980 0.0333655777
## [29,] 7985 0.0351301920
## [30,] 7990 0.0349099754
## [31,] 7995 0.0328448889
## [32,] 8000 0.0316308556
## [33,] 8005 0.0338190861
## [34,] 8010 0.0389705163
## [35,] 8015 0.0446963589
## [36,] 8020 0.0457789840
## [37,] 8025 0.0393994960
## [38,] 8030 0.0306946215
## [39,] 8035 0.0234615710
## [40,] 8040 0.0179094234
## [41,] 8045 0.0146274566
## [42,] 8050 0.0124116660
## [43,] 8055 0.0110048596
## [44,] 8060 0.0104816198
## [45,] 8065 0.0106199687
## [46,] 8070 0.0111408561
## [47,] 8075 0.0115504348
## [48,] 8080 0.0118668677
## [49,] 8085 0.0128714983
## [50,] 8090 0.0150275241
## [51,] 8095 0.0169858658
## [52,] 8100 0.0190578652
## [53,] 8105 0.0194752106
## [54,] 8110 0.0190102962
## [55,] 8115 0.0172536612
## [56,] 8120 0.0156836737
## [57,] 8125 0.0149396568
## [58,] 8130 0.0144907542
## [59,] 8135 0.0156093128
## [60,] 8140 0.0167914021
## [61,] 8145 0.0184650320
## [62,] 8150 0.0197911617
## [63,] 8155 0.0202461925
## [64,] 8160 0.0189310091
## [65,] 8165 0.0150275241
## [66,] 8170 0.0094437433
## [67,] 8175 0.0052252879
## [68,] 8180 0.0027137563
## [69,] 8185 0.0015263656
## [70,] 8190 0.0010331326
## [71,] 8195 0.0006939167
## [72,] 8200 0.0005844432
## [73,] 8205 0.0005054591
## [74,] 8210 0.0004379560
## [75,] 8215 0.0003713799
## [76,] 8220 0.0002710205
## [77,] 8225 0.0002735940
## [78,] 8230 0.0002559174
## [79,] 8235 0.0002309766
## [80,] 8240 0.0002710205
## [81,] 8245 0.0003042909
## [82,] 8250 0.0003003616
## [83,] 8255 0.0003225044
## [84,] 8260 0.0003723607
## [85,] 8265 0.0004993458
## [86,] 8270 0.0006579833
## [87,] 8275 0.0007661465
## [88,] 8280 0.0006098987
## [89,] 8285 0.0005145501
## 
## $calib$probs[[4]]
##         [,1]         [,2]
##   [1,]  9425 0.0004865336
##   [2,]  9430 0.0006647285
##   [3,]  9435 0.0007973484
##   [4,]  9440 0.0008831414
##   [5,]  9445 0.0009393424
##   [6,]  9450 0.0010504317
##   [7,]  9455 0.0011088173
##   [8,]  9460 0.0011420989
##   [9,]  9465 0.0013227820
##  [10,]  9470 0.0016468369
##  [11,]  9475 0.0022927852
##  [12,]  9480 0.0032570683
##  [13,]  9485 0.0044420705
##  [14,]  9490 0.0059908243
##  [15,]  9495 0.0080431303
##  [16,]  9500 0.0087132134
##  [17,]  9505 0.0083417200
##  [18,]  9510 0.0066943523
##  [19,]  9515 0.0054214901
##  [20,]  9520 0.0054214901
##  [21,]  9525 0.0073735341
##  [22,]  9530 0.0125427544
##  [23,]  9535 0.0224147451
##  [24,]  9540 0.0387473468
##  [25,]  9545 0.0460995880
##  [26,]  9550 0.0392531537
##  [27,]  9555 0.0293680192
##  [28,]  9560 0.0220211657
##  [29,]  9565 0.0197415586
##  [30,]  9570 0.0216313221
##  [31,]  9575 0.0255552973
##  [32,]  9580 0.0280833270
##  [33,]  9585 0.0296295441
##  [34,]  9590 0.0287800131
##  [35,]  9595 0.0275134519
##  [36,]  9600 0.0243221601
##  [37,]  9605 0.0211724372
##  [38,]  9610 0.0196401661
##  [39,]  9615 0.0200010314
##  [40,]  9620 0.0211075736
##  [41,]  9625 0.0203659362
##  [42,]  9630 0.0185824347
##  [43,]  9635 0.0178985323
##  [44,]  9640 0.0186166453
##  [45,]  9645 0.0207950694
##  [46,]  9650 0.0227196811
##  [47,]  9655 0.0219387012
##  [48,]  9660 0.0193760877
##  [49,]  9665 0.0179247239
##  [50,]  9670 0.0175854046
##  [51,]  9675 0.0178790491
##  [52,]  9680 0.0168978188
##  [53,]  9685 0.0136626960
##  [54,]  9690 0.0105031353
##  [55,]  9695 0.0082805573
##  [56,]  9700 0.0069059209
##  [57,]  9705 0.0057993231
##  [58,]  9710 0.0052428348
##  [59,]  9715 0.0051808059
##  [60,]  9720 0.0049518562
##  [61,]  9725 0.0047328574
##  [62,]  9730 0.0046269687
##  [63,]  9735 0.0045819694
##  [64,]  9740 0.0044804755
##  [65,]  9745 0.0042263891
##  [66,]  9750 0.0039849589
##  [67,]  9755 0.0038948219
##  [68,]  9760 0.0036365164
##  [69,]  9765 0.0033954844
##  [70,]  9770 0.0030991005
##  [71,]  9775 0.0026580541
##  [72,]  9780 0.0024684912
##  [73,]  9785 0.0022547822
##  [74,]  9790 0.0021978856
##  [75,]  9795 0.0021978856
##  [76,]  9800 0.0022927852
##  [77,]  9805 0.0023400055
##  [78,]  9810 0.0024416730
##  [79,]  9815 0.0021455625
##  [80,]  9820 0.0019228246
##  [81,]  9825 0.0018393387
##  [82,]  9830 0.0017596449
##  [83,]  9835 0.0018393387
##  [84,]  9840 0.0020555617
##  [85,]  9845 0.0023504844
##  [86,]  9850 0.0026784563
##  [87,]  9855 0.0026296594
##  [88,]  9860 0.0028636589
##  [89,]  9865 0.0031857886
##  [90,]  9870 0.0032949049
##  [91,]  9875 0.0036069274
##  [92,]  9880 0.0034473426
##  [93,]  9885 0.0027032041
##  [94,]  9890 0.0021019047
##  [95,]  9895 0.0015081424
##  [96,]  9900 0.0011179987
##  [97,]  9905 0.0007751046
##  [98,]  9910 0.0005919267
##  [99,]  9915 0.0005093586
## [100,]  9920 0.0004677294
## [101,]  9925 0.0004409924
## [102,]  9930 0.0004657507
## [103,]  9935 0.0005114011
## [104,]  9940 0.0005528378
## [105,]  9945 0.0005667551
## [106,]  9950 0.0005896854
## [107,]  9955 0.0005163016
## [108,]  9960 0.0003802887
## [109,]  9965 0.0002586455
## [110,]  9970 0.0001842331
## [111,]  9975 0.0001619774
## [112,]  9980 0.0001911983
## [113,]  9985 0.0002734568
## [114,]  9990 0.0004015426
## [115,]  9995 0.0004830342
## [116,] 10000 0.0005237278
## [117,] 10005 0.0005480622
## [118,] 10010 0.0004095000
## [119,] 10015 0.0002952405
## [120,] 10020 0.0003009798
## [121,] 10025 0.0003786676
## [122,] 10030 0.0004586256
## [123,] 10035 0.0004259368
## [124,] 10040 0.0003642116
## [125,] 10045 0.0002632329
## [126,] 10050 0.0002060188
## [127,] 10055 0.0002221525
## [128,] 10060 0.0002896212
## [129,] 10065 0.0004198678
## [130,] 10070 0.0005747823
## [131,] 10075 0.0006295921
## [132,] 10080 0.0006765964
## [133,] 10085 0.0006427182
## [134,] 10090 0.0006107047
## [135,] 10095 0.0006312910
## [136,] 10100 0.0007373563
## [137,] 10105 0.0007982036
## [138,] 10110 0.0006628627
## [139,] 10115 0.0005093586
## 
## $calib$probs[[5]]
##         [,1]         [,2]
##   [1,] 10710 0.0003487031
##   [2,] 10715 0.0004202430
##   [3,] 10720 0.0005676822
##   [4,] 10725 0.0007422027
##   [5,] 10730 0.0008838234
##   [6,] 10735 0.0009442099
##   [7,] 10740 0.0009911761
##   [8,] 10745 0.0010428037
##   [9,] 10750 0.0010333016
##  [10,] 10755 0.0010518453
##  [11,] 10760 0.0011876103
##  [12,] 10765 0.0012952676
##  [13,] 10770 0.0015421006
##  [14,] 10775 0.0017374533
##  [15,] 10780 0.0019671386
##  [16,] 10785 0.0024221859
##  [17,] 10790 0.0030547683
##  [18,] 10795 0.0036568203
##  [19,] 10800 0.0039745232
##  [20,] 10805 0.0045244454
##  [21,] 10810 0.0052262224
##  [22,] 10815 0.0060086327
##  [23,] 10820 0.0068740597
##  [24,] 10825 0.0069600278
##  [25,] 10830 0.0070533444
##  [26,] 10835 0.0070533444
##  [27,] 10840 0.0063105260
##  [28,] 10845 0.0050694165
##  [29,] 10850 0.0041106824
##  [30,] 10855 0.0034896777
##  [31,] 10860 0.0032271631
##  [32,] 10865 0.0038492964
##  [33,] 10870 0.0053599505
##  [34,] 10875 0.0077787838
##  [35,] 10880 0.0098211361
##  [36,] 10885 0.0110038493
##  [37,] 10890 0.0120957773
##  [38,] 10895 0.0132204534
##  [39,] 10900 0.0147100402
##  [40,] 10905 0.0153191228
##  [41,] 10910 0.0150128169
##  [42,] 10915 0.0147121279
##  [43,] 10920 0.0132739007
##  [44,] 10925 0.0111259508
##  [45,] 10930 0.0091421745
##  [46,] 10935 0.0083810145
##  [47,] 10940 0.0083810145
##  [48,] 10945 0.0072645790
##  [49,] 10950 0.0054806826
##  [50,] 10955 0.0046872926
##  [51,] 10960 0.0040790773
##  [52,] 10965 0.0032271631
##  [53,] 10970 0.0022401325
##  [54,] 10975 0.0014707212
##  [55,] 10980 0.0011468742
##  [56,] 10985 0.0010364632
##  [57,] 10990 0.0011763717
##  [58,] 10995 0.0015994314
##  [59,] 11000 0.0022992176
##  [60,] 11005 0.0035641446
##  [61,] 11010 0.0048435742
##  [62,] 11015 0.0056370141
##  [63,] 11020 0.0052803198
##  [64,] 11025 0.0044089162
##  [65,] 11030 0.0031567590
##  [66,] 11035 0.0023597009
##  [67,] 11040 0.0019109344
##  [68,] 11045 0.0017204115
##  [69,] 11050 0.0015095630
##  [70,] 11055 0.0015191544
##  [71,] 11060 0.0017661240
##  [72,] 11065 0.0027742140
##  [73,] 11070 0.0055506458
##  [74,] 11075 0.0097709042
##  [75,] 11080 0.0126495445
##  [76,] 11085 0.0138081236
##  [77,] 11090 0.0141080945
##  [78,] 11095 0.0147327721
##  [79,] 11100 0.0173115674
##  [80,] 11105 0.0200797817
##  [81,] 11110 0.0229212575
##  [82,] 11115 0.0246733026
##  [83,] 11120 0.0248633200
##  [84,] 11125 0.0247162888
##  [85,] 11130 0.0240412958
##  [86,] 11135 0.0234894704
##  [87,] 11140 0.0239780781
##  [88,] 11145 0.0251779169
##  [89,] 11150 0.0265313510
##  [90,] 11155 0.0276304948
##  [91,] 11160 0.0284625255
##  [92,] 11165 0.0299469450
##  [93,] 11170 0.0315917966
##  [94,] 11175 0.0327780711
##  [95,] 11180 0.0313828541
##  [96,] 11185 0.0289872853
##  [97,] 11190 0.0258625854
##  [98,] 11195 0.0205343616
##  [99,] 11200 0.0147199440
## [100,] 11205 0.0105859890
## [101,] 11210 0.0077076729
## [102,] 11215 0.0065490985
## [103,] 11220 0.0058620653
## [104,] 11225 0.0049680909
## [105,] 11230 0.0042642256
## [106,] 11235 0.0030631140
## [107,] 11240 0.0020252646
## [108,] 11245 0.0012952676
## [109,] 11250 0.0007997823
## [110,] 11255 0.0005612079
## [111,] 11260 0.0004618648
## [112,] 11265 0.0003387672
## 
## $calib$probs[[6]]
##         [,1]         [,2]
##   [1,] 11240 1.665541e-04
##   [2,] 11245 2.358136e-04
##   [3,] 11250 3.006180e-04
##   [4,] 11255 3.284416e-04
##   [5,] 11260 3.754758e-04
##   [6,] 11265 4.678720e-04
##   [7,] 11270 6.051985e-04
##   [8,] 11275 7.778892e-04
##   [9,] 11280 8.784969e-04
##  [10,] 11285 9.488030e-04
##  [11,] 11290 9.728322e-04
##  [12,] 11295 9.728322e-04
##  [13,] 11300 9.977221e-04
##  [14,] 11305 1.041043e-03
##  [15,] 11310 1.123607e-03
##  [16,] 11315 1.213548e-03
##  [17,] 11320 1.358670e-03
##  [18,] 11325 1.558618e-03
##  [19,] 11330 1.727366e-03
##  [20,] 11335 1.849637e-03
##  [21,] 11340 2.099318e-03
##  [22,] 11345 2.447358e-03
##  [23,] 11350 2.883390e-03
##  [24,] 11355 3.331617e-03
##  [25,] 11360 3.726873e-03
##  [26,] 11365 3.787628e-03
##  [27,] 11370 3.384469e-03
##  [28,] 11375 2.679492e-03
##  [29,] 11380 2.228612e-03
##  [30,] 11385 2.133860e-03
##  [31,] 11390 1.906214e-03
##  [32,] 11395 1.608151e-03
##  [33,] 11400 1.471333e-03
##  [34,] 11405 1.670182e-03
##  [35,] 11410 1.955849e-03
##  [36,] 11415 2.215341e-03
##  [37,] 11420 2.432064e-03
##  [38,] 11425 2.872501e-03
##  [39,] 11430 3.413372e-03
##  [40,] 11435 3.445165e-03
##  [41,] 11440 2.921842e-03
##  [42,] 11445 2.316988e-03
##  [43,] 11450 1.876945e-03
##  [44,] 11455 1.933742e-03
##  [45,] 11460 2.149287e-03
##  [46,] 11465 2.174239e-03
##  [47,] 11470 1.979378e-03
##  [48,] 11475 1.745776e-03
##  [49,] 11480 1.568378e-03
##  [50,] 11485 1.453063e-03
##  [51,] 11490 1.312539e-03
##  [52,] 11495 1.239680e-03
##  [53,] 11500 1.305518e-03
##  [54,] 11505 1.435350e-03
##  [55,] 11510 1.734552e-03
##  [56,] 11515 2.415143e-03
##  [57,] 11520 3.413372e-03
##  [58,] 11525 4.133422e-03
##  [59,] 11530 4.162310e-03
##  [60,] 11535 3.946572e-03
##  [61,] 11540 3.996093e-03
##  [62,] 11545 4.352934e-03
##  [63,] 11550 4.171903e-03
##  [64,] 11555 3.557181e-03
##  [65,] 11560 2.827075e-03
##  [66,] 11565 2.387190e-03
##  [67,] 11570 2.312732e-03
##  [68,] 11575 2.226941e-03
##  [69,] 11580 2.064313e-03
##  [70,] 11585 1.758413e-03
##  [71,] 11590 1.372129e-03
##  [72,] 11595 1.127452e-03
##  [73,] 11600 1.105994e-03
##  [74,] 11605 1.279952e-03
##  [75,] 11610 1.598590e-03
##  [76,] 11615 1.949843e-03
##  [77,] 11620 2.496921e-03
##  [78,] 11625 3.179809e-03
##  [79,] 11630 4.117342e-03
##  [80,] 11635 5.219007e-03
##  [81,] 11640 6.723955e-03
##  [82,] 11645 8.152103e-03
##  [83,] 11650 8.977858e-03
##  [84,] 11655 8.890152e-03
##  [85,] 11660 7.780628e-03
##  [86,] 11665 6.441768e-03
##  [87,] 11670 5.401773e-03
##  [88,] 11675 4.688494e-03
##  [89,] 11680 4.644543e-03
##  [90,] 11685 4.970486e-03
##  [91,] 11690 5.189310e-03
##  [92,] 11695 4.851153e-03
##  [93,] 11700 4.311454e-03
##  [94,] 11705 4.543206e-03
##  [95,] 11710 5.821953e-03
##  [96,] 11715 7.780628e-03
##  [97,] 11720 9.568856e-03
##  [98,] 11725 1.035511e-02
##  [99,] 11730 1.059287e-02
## [100,] 11735 1.081075e-02
## [101,] 11740 1.082928e-02
## [102,] 11745 9.900499e-03
## [103,] 11750 8.715446e-03
## [104,] 11755 9.124526e-03
## [105,] 11760 1.128576e-02
## [106,] 11765 1.303089e-02
## [107,] 11770 1.346534e-02
## [108,] 11775 1.370831e-02
## [109,] 11780 1.427607e-02
## [110,] 11785 1.473544e-02
## [111,] 11790 1.487871e-02
## [112,] 11795 1.479378e-02
## [113,] 11800 1.428767e-02
## [114,] 11805 1.372134e-02
## [115,] 11810 1.355821e-02
## [116,] 11815 1.355058e-02
## [117,] 11820 1.310430e-02
## [118,] 11825 1.270310e-02
## [119,] 11830 1.310430e-02
## [120,] 11835 1.428357e-02
## [121,] 11840 1.489123e-02
## [122,] 11845 1.454385e-02
## [123,] 11850 1.421103e-02
## [124,] 11855 1.427607e-02
## [125,] 11860 1.454385e-02
## [126,] 11865 1.466824e-02
## [127,] 11870 1.459018e-02
## [128,] 11875 1.396116e-02
## [129,] 11880 1.265217e-02
## [130,] 11885 1.119893e-02
## [131,] 11890 1.061367e-02
## [132,] 11895 1.063591e-02
## [133,] 11900 1.075344e-02
## [134,] 11905 1.135431e-02
## [135,] 11910 1.204138e-02
## [136,] 11915 1.244772e-02
## [137,] 11920 1.213103e-02
## [138,] 11925 1.103277e-02
## [139,] 11930 9.421507e-03
## [140,] 11935 8.004192e-03
## [141,] 11940 7.127590e-03
## [142,] 11945 6.673606e-03
## [143,] 11950 6.673606e-03
## [144,] 11955 6.571128e-03
## [145,] 11960 6.073278e-03
## [146,] 11965 5.419065e-03
## [147,] 11970 4.654541e-03
## [148,] 11975 4.089333e-03
## [149,] 11980 3.946542e-03
## [150,] 11985 4.089333e-03
## [151,] 11990 4.120038e-03
## [152,] 11995 3.837820e-03
## [153,] 12000 3.475577e-03
## [154,] 12005 3.383207e-03
## [155,] 12010 3.508105e-03
## [156,] 12015 3.262204e-03
## [157,] 12020 2.663013e-03
## [158,] 12025 2.356783e-03
## [159,] 12030 2.517975e-03
## [160,] 12035 2.738354e-03
## [161,] 12040 2.743407e-03
## [162,] 12045 2.743407e-03
## [163,] 12050 2.842233e-03
## [164,] 12055 2.763928e-03
## [165,] 12060 2.424858e-03
## [166,] 12065 2.144853e-03
## [167,] 12070 2.144163e-03
## [168,] 12075 2.335977e-03
## [169,] 12080 2.517975e-03
## [170,] 12085 2.405384e-03
## [171,] 12090 1.906214e-03
## [172,] 12095 1.305518e-03
## [173,] 12100 9.322165e-04
## [174,] 12105 8.201802e-04
## [175,] 12110 9.500578e-04
## [176,] 12115 1.216308e-03
## [177,] 12120 1.509633e-03
## [178,] 12125 1.706053e-03
## [179,] 12130 1.826615e-03
## [180,] 12135 1.841489e-03
## [181,] 12140 1.792199e-03
## [182,] 12145 1.557735e-03
## [183,] 12150 1.271165e-03
## [184,] 12155 1.155374e-03
## [185,] 12160 1.247109e-03
## [186,] 12165 1.247109e-03
## [187,] 12170 1.064681e-03
## [188,] 12175 8.047853e-04
## [189,] 12180 6.789958e-04
## [190,] 12185 6.207825e-04
## [191,] 12190 6.207825e-04
## [192,] 12195 5.867611e-04
## [193,] 12200 5.129201e-04
## [194,] 12205 4.178969e-04
## [195,] 12210 3.552412e-04
## [196,] 12215 3.884982e-04
## [197,] 12220 5.211447e-04
## [198,] 12225 6.821453e-04
## [199,] 12230 7.937906e-04
## [200,] 12235 8.518883e-04
## [201,] 12240 9.147183e-04
## [202,] 12245 1.064681e-03
## [203,] 12250 1.193378e-03
## [204,] 12255 1.127173e-03
## [205,] 12260 9.284814e-04
## [206,] 12265 7.461328e-04
## [207,] 12270 7.050636e-04
## [208,] 12275 6.639461e-04
## [209,] 12280 5.509565e-04
## [210,] 12285 4.665084e-04
## [211,] 12290 4.186607e-04
## [212,] 12295 4.264178e-04
## [213,] 12300 5.022180e-04
## [214,] 12305 6.160856e-04
## [215,] 12310 7.603680e-04
## [216,] 12315 8.730245e-04
## [217,] 12320 8.642190e-04
## [218,] 12325 6.860820e-04
## [219,] 12330 4.657906e-04
## [220,] 12335 3.364468e-04
## [221,] 12340 3.251359e-04
## [222,] 12345 4.035929e-04
## [223,] 12350 5.706453e-04
## [224,] 12355 7.292800e-04
## [225,] 12360 8.848603e-04
## [226,] 12365 9.978345e-04
## [227,] 12370 1.077195e-03
## [228,] 12375 1.070450e-03
## [229,] 12380 9.191452e-04
## [230,] 12385 8.090215e-04
## [231,] 12390 8.682132e-04
## [232,] 12395 1.010937e-03
## [233,] 12400 1.200359e-03
## [234,] 12405 1.329989e-03
## [235,] 12410 1.453063e-03
## [236,] 12415 1.520968e-03
## [237,] 12420 1.580358e-03
## [238,] 12425 1.453063e-03
## [239,] 12430 1.165141e-03
## [240,] 12435 8.449515e-04
## [241,] 12440 6.539307e-04
## [242,] 12445 5.978791e-04
## [243,] 12450 6.092181e-04
## [244,] 12455 6.092181e-04
## [245,] 12460 5.651707e-04
## [246,] 12465 4.590048e-04
## [247,] 12470 3.422109e-04
## [248,] 12475 2.485649e-04
## [249,] 12480 1.851025e-04
## [250,] 12485 1.343552e-04
## [251,] 12490 8.430422e-05
## [252,] 12495 5.223501e-05
## [253,] 12500 4.249777e-05
## [254,] 12505 4.124649e-05
## [255,] 12510 3.758667e-05
## [256,] 12515 3.306675e-05
## [257,] 12520 2.821504e-05
## [258,] 12525 2.920585e-05
## [259,] 12530 4.324709e-05
## [260,] 12535 7.256679e-05
## [261,] 12540 1.061923e-04
## [262,] 12545 1.268125e-04
## [263,] 12550 1.333599e-04
## [264,] 12555 1.379447e-04
## [265,] 12560 1.579937e-04
## [266,] 12565 1.722810e-04
## [267,] 12570 1.721360e-04
## [268,] 12575 1.607176e-04
## 
## $calib$probs[[7]]
##        [,1]         [,2]
##  [1,] 12725 0.0005223418
##  [2,] 12730 0.0010184326
##  [3,] 12735 0.0021247793
##  [4,] 12740 0.0045813584
##  [5,] 12745 0.0098774529
##  [6,] 12750 0.0184894386
##  [7,] 12755 0.0254225359
##  [8,] 12760 0.0292978198
##  [9,] 12765 0.0340512343
## [10,] 12770 0.0398492351
## [11,] 12775 0.0402394328
## [12,] 12780 0.0389808147
## [13,] 12785 0.0375821297
## [14,] 12790 0.0338754013
## [15,] 12795 0.0306109433
## [16,] 12800 0.0306109433
## [17,] 12805 0.0332781694
## [18,] 12810 0.0358297640
## [19,] 12815 0.0382188060
## [20,] 12820 0.0396249800
## [21,] 12825 0.0396249800
## [22,] 12830 0.0376207454
## [23,] 12835 0.0334452873
## [24,] 12840 0.0279930407
## [25,] 12845 0.0218064056
## [26,] 12850 0.0166407934
## [27,] 12855 0.0138319381
## [28,] 12860 0.0129708929
## [29,] 12865 0.0138319381
## [30,] 12870 0.0156702308
## [31,] 12875 0.0195627589
## [32,] 12880 0.0245122133
## [33,] 12885 0.0267686905
## [34,] 12890 0.0246543969
## [35,] 12895 0.0189558948
## [36,] 12900 0.0132725120
## [37,] 12905 0.0109135399
## [38,] 12910 0.0108326747
## [39,] 12915 0.0102226354
## [40,] 12920 0.0078214234
## [41,] 12925 0.0048213006
## [42,] 12930 0.0030941227
## [43,] 12935 0.0023988251
## [44,] 12940 0.0024662642
## [45,] 12945 0.0030098710
## [46,] 12950 0.0038742132
## [47,] 12955 0.0050914073
## [48,] 12960 0.0058389355
## [49,] 12965 0.0050849719
## [50,] 12970 0.0031984991
## [51,] 12975 0.0017395384
## [52,] 12980 0.0012441164
## [53,] 12985 0.0014015210
## [54,] 12990 0.0016768445
## [55,] 12995 0.0018073389
## [56,] 13000 0.0016534401
## [57,] 13005 0.0014396965
## [58,] 13010 0.0014632380
## [59,] 13015 0.0013859885
## [60,] 13020 0.0011984494
## [61,] 13025 0.0010472760
## [62,] 13030 0.0010045512
## [63,] 13035 0.0011937968
## [64,] 13040 0.0016665579
## [65,] 13045 0.0018711548
## [66,] 13050 0.0017000473
## [67,] 13055 0.0015140356
## [68,] 13060 0.0014320605
## [69,] 13065 0.0013177556
## [70,] 13070 0.0011353681
## [71,] 13075 0.0008461378
## [72,] 13080 0.0006553321
## [73,] 13085 0.0004375149
## 
## $calib$probs[[8]]
##        [,1]         [,2]
##  [1,] 12485 0.0011057559
##  [2,] 12490 0.0021057724
##  [3,] 12495 0.0033727320
##  [4,] 12500 0.0042016379
##  [5,] 12505 0.0044101581
##  [6,] 12510 0.0048852582
##  [7,] 12515 0.0063467198
##  [8,] 12520 0.0082827764
##  [9,] 12525 0.0082259023
## [10,] 12530 0.0058285460
## [11,] 12535 0.0030108536
## [12,] 12540 0.0014239010
## [13,] 12545 0.0008423176
## [14,] 12550 0.0007514571
## [15,] 12555 0.0007192297
## [16,] 12560 0.0006345339
## [17,] 12565 0.0005431467
## [18,] 12570 0.0005701047
## [19,] 12575 0.0006210382
## [20,] 12580 0.0007225285
## [21,] 12585 0.0008423176
## [22,] 12590 0.0010569849
## [23,] 12595 0.0012654479
## [24,] 12600 0.0016271755
## [25,] 12605 0.0023189774
## [26,] 12610 0.0035453650
## [27,] 12615 0.0062519510
## [28,] 12620 0.0104517200
## [29,] 12625 0.0146840242
## [30,] 12630 0.0179682658
## [31,] 12635 0.0213681305
## [32,] 12640 0.0249890587
## [33,] 12645 0.0255031377
## [34,] 12650 0.0222903672
## [35,] 12655 0.0200390524
## [36,] 12660 0.0204650666
## [37,] 12665 0.0213684632
## [38,] 12670 0.0202431142
## [39,] 12675 0.0171899176
## [40,] 12680 0.0151564782
## [41,] 12685 0.0182035606
## [42,] 12690 0.0281179383
## [43,] 12695 0.0425146375
## [44,] 12700 0.0536576782
## [45,] 12705 0.0592496157
## [46,] 12710 0.0609913553
## [47,] 12715 0.0609913553
## [48,] 12720 0.0609913553
## [49,] 12725 0.0631199041
## [50,] 12730 0.0608454655
## [51,] 12735 0.0484423723
## [52,] 12740 0.0322943965
## [53,] 12745 0.0181882172
## [54,] 12750 0.0091629229
## [55,] 12755 0.0056555896
## [56,] 12760 0.0046527262
## [57,] 12765 0.0038043114
## [58,] 12770 0.0024802502
## [59,] 12775 0.0016229401
## [60,] 12780 0.0013488410
## [61,] 12785 0.0012021817
## [62,] 12790 0.0009809349
## [63,] 12795 0.0008189520
## [64,] 12800 0.0008189520
## [65,] 12805 0.0009158757
## [66,] 12810 0.0010725242
## [67,] 12815 0.0012879800
## [68,] 12820 0.0014459567
## [69,] 12825 0.0014459567
## [70,] 12830 0.0012303448
## [71,] 12835 0.0009591738
## [72,] 12840 0.0007547442
## [73,] 12845 0.0005701047
## [74,] 12850 0.0004428103
## [75,] 12855 0.0003725717
## [76,] 12860 0.0003482588
## [77,] 12865 0.0003725717
## [78,] 12870 0.0004131278
## [79,] 12875 0.0005127562
## [80,] 12880 0.0006726334
## [81,] 12885 0.0007809883
## [82,] 12890 0.0007264531
## 
## $calib$probs[[9]]
##         [,1]         [,2]
##   [1,] 12095 0.0004291513
##   [2,] 12100 0.0006449079
##   [3,] 12105 0.0007290868
##   [4,] 12110 0.0006276185
##   [5,] 12115 0.0004432025
##   [6,] 12120 0.0003098387
##   [7,] 12125 0.0002478200
##   [8,] 12130 0.0002402423
##   [9,] 12135 0.0002239770
##  [10,] 12140 0.0002463632
##  [11,] 12145 0.0003161906
##  [12,] 12150 0.0003916492
##  [13,] 12155 0.0004470726
##  [14,] 12160 0.0004021050
##  [15,] 12165 0.0004021050
##  [16,] 12170 0.0005338478
##  [17,] 12175 0.0007495881
##  [18,] 12180 0.0009645016
##  [19,] 12185 0.0010405543
##  [20,] 12190 0.0010405543
##  [21,] 12195 0.0011341425
##  [22,] 12200 0.0014817217
##  [23,] 12205 0.0021770426
##  [24,] 12210 0.0026624985
##  [25,] 12215 0.0024467915
##  [26,] 12220 0.0015376991
##  [27,] 12225 0.0009023612
##  [28,] 12230 0.0007204645
##  [29,] 12235 0.0006898696
##  [30,] 12240 0.0006627118
##  [31,] 12245 0.0005338478
##  [32,] 12250 0.0004550601
##  [33,] 12255 0.0004927490
##  [34,] 12260 0.0006925321
##  [35,] 12265 0.0008379958
##  [36,] 12270 0.0009116085
##  [37,] 12275 0.0010618563
##  [38,] 12280 0.0014108328
##  [39,] 12285 0.0018285731
##  [40,] 12290 0.0020423205
##  [41,] 12295 0.0019832702
##  [42,] 12300 0.0016288760
##  [43,] 12305 0.0011890163
##  [44,] 12310 0.0008148931
##  [45,] 12315 0.0006270837
##  [46,] 12320 0.0007193741
##  [47,] 12325 0.0011543945
##  [48,] 12330 0.0020880126
##  [49,] 12335 0.0029095125
##  [50,] 12340 0.0029072774
##  [51,] 12345 0.0021659321
##  [52,] 12350 0.0014260179
##  [53,] 12355 0.0009228902
##  [54,] 12360 0.0006529134
##  [55,] 12365 0.0005176333
##  [56,] 12370 0.0004645244
##  [57,] 12375 0.0004975596
##  [58,] 12380 0.0006180818
##  [59,] 12385 0.0007006544
##  [60,] 12390 0.0006711176
##  [61,] 12395 0.0005394783
##  [62,] 12400 0.0004239406
##  [63,] 12405 0.0003467131
##  [64,] 12410 0.0003262566
##  [65,] 12415 0.0002889772
##  [66,] 12420 0.0002744799
##  [67,] 12425 0.0003262566
##  [68,] 12430 0.0005023126
##  [69,] 12435 0.0007935857
##  [70,] 12440 0.0010206783
##  [71,] 12445 0.0011019908
##  [72,] 12450 0.0010708047
##  [73,] 12455 0.0010708047
##  [74,] 12460 0.0012014660
##  [75,] 12465 0.0017642082
##  [76,] 12470 0.0031963616
##  [77,] 12475 0.0053867605
##  [78,] 12480 0.0087126048
##  [79,] 12485 0.0166022742
##  [80,] 12490 0.0290230281
##  [81,] 12495 0.0370394241
##  [82,] 12500 0.0379497161
##  [83,] 12505 0.0378724648
##  [84,] 12510 0.0373371146
##  [85,] 12515 0.0349849738
##  [86,] 12520 0.0308312858
##  [87,] 12525 0.0313552896
##  [88,] 12530 0.0372673586
##  [89,] 12535 0.0342082856
##  [90,] 12540 0.0217316949
##  [91,] 12545 0.0137491486
##  [92,] 12550 0.0121889664
##  [93,] 12555 0.0115456940
##  [94,] 12560 0.0096787371
##  [95,] 12565 0.0080083243
##  [96,] 12570 0.0084050068
##  [97,] 12575 0.0094115352
##  [98,] 12580 0.0114178130
##  [99,] 12585 0.0137491486
## [100,] 12590 0.0174217622
## [101,] 12595 0.0208810309
## [102,] 12600 0.0259701285
## [103,] 12605 0.0319636852
## [104,] 12610 0.0370728185
## [105,] 12615 0.0367402944
## [106,] 12620 0.0288681559
## [107,] 12625 0.0206366350
## [108,] 12630 0.0168234790
## [109,] 12635 0.0138115903
## [110,] 12640 0.0110035388
## [111,] 12645 0.0107087072
## [112,] 12650 0.0131213927
## [113,] 12655 0.0148948034
## [114,] 12660 0.0137491486
## [115,] 12665 0.0130495555
## [116,] 12670 0.0143128160
## [117,] 12675 0.0176356455
## [118,] 12680 0.0195490329
## [119,] 12685 0.0170480920
## [120,] 12690 0.0107033279
## [121,] 12695 0.0050850730
## [122,] 12700 0.0027423792
## [123,] 12705 0.0018032109
## [124,] 12710 0.0015071390
## [125,] 12715 0.0015071390
## [126,] 12720 0.0015071390
## [127,] 12725 0.0012463153
## [128,] 12730 0.0007859940
## [129,] 12735 0.0004164305
## 
## $calib$probs[[10]]
##         [,1]         [,2]
##   [1,] 12750 0.0003166976
##   [2,] 12755 0.0004080680
##   [3,] 12760 0.0004854824
##   [4,] 12765 0.0006338990
##   [5,] 12770 0.0009926599
##   [6,] 12775 0.0013117619
##   [7,] 12780 0.0015166397
##   [8,] 12785 0.0017055175
##   [9,] 12790 0.0021943132
##  [10,] 12795 0.0025468067
##  [11,] 12800 0.0025468067
##  [12,] 12805 0.0021771388
##  [13,] 12810 0.0019185728
##  [14,] 12815 0.0016559955
##  [15,] 12820 0.0014137946
##  [16,] 12825 0.0014137946
##  [17,] 12830 0.0017353441
##  [18,] 12835 0.0022465009
##  [19,] 12840 0.0030260954
##  [20,] 12845 0.0040933261
##  [21,] 12850 0.0053860280
##  [22,] 12855 0.0062280842
##  [23,] 12860 0.0064425936
##  [24,] 12865 0.0062280842
##  [25,] 12870 0.0055653223
##  [26,] 12875 0.0045923641
##  [27,] 12880 0.0037656104
##  [28,] 12885 0.0035628640
##  [29,] 12890 0.0040651604
##  [30,] 12895 0.0054268096
##  [31,] 12900 0.0071415558
##  [32,] 12905 0.0081863001
##  [33,] 12910 0.0084567419
##  [34,] 12915 0.0091002331
##  [35,] 12920 0.0117465230
##  [36,] 12925 0.0156599744
##  [37,] 12930 0.0194390628
##  [38,] 12935 0.0214024228
##  [39,] 12940 0.0211536073
##  [40,] 12945 0.0196970286
##  [41,] 12950 0.0177102501
##  [42,] 12955 0.0151578270
##  [43,] 12960 0.0132394698
##  [44,] 12965 0.0144548419
##  [45,] 12970 0.0195183590
##  [46,] 12975 0.0243811810
##  [47,] 12980 0.0261978165
##  [48,] 12985 0.0257859929
##  [49,] 12990 0.0243477078
##  [50,] 12995 0.0234708261
##  [51,] 13000 0.0238676866
##  [52,] 13005 0.0248711971
##  [53,] 13010 0.0252790284
##  [54,] 13015 0.0256110329
##  [55,] 13020 0.0261847454
##  [56,] 13025 0.0267670850
##  [57,] 13030 0.0270291039
##  [58,] 13035 0.0265747295
##  [59,] 13040 0.0249790176
##  [60,] 13045 0.0235110938
##  [61,] 13050 0.0236521795
##  [62,] 13055 0.0242550041
##  [63,] 13060 0.0246550653
##  [64,] 13065 0.0252133329
##  [65,] 13070 0.0264397491
##  [66,] 13075 0.0272950570
##  [67,] 13080 0.0270330563
##  [68,] 13085 0.0246087383
##  [69,] 13090 0.0194315594
##  [70,] 13095 0.0128662820
##  [71,] 13100 0.0082114563
##  [72,] 13105 0.0056604590
##  [73,] 13110 0.0045935716
##  [74,] 13115 0.0037153492
##  [75,] 13120 0.0028445907
##  [76,] 13125 0.0026549164
##  [77,] 13130 0.0031078216
##  [78,] 13135 0.0038138509
##  [79,] 13140 0.0045462399
##  [80,] 13145 0.0057166641
##  [81,] 13150 0.0062883196
##  [82,] 13155 0.0055611731
##  [83,] 13160 0.0041418342
##  [84,] 13165 0.0026951842
##  [85,] 13170 0.0017294759
##  [86,] 13175 0.0011370764
##  [87,] 13180 0.0007523862
##  [88,] 13185 0.0005178802
##  [89,] 13190 0.0004192051
##  [90,] 13195 0.0004056213
##  [91,] 13200 0.0004212215
##  [92,] 13205 0.0004720401
##  [93,] 13210 0.0005398197
##  [94,] 13215 0.0006181607
##  [95,] 13220 0.0006959712
##  [96,] 13225 0.0006694697
##  [97,] 13230 0.0005628728
##  [98,] 13235 0.0004508795
##  [99,] 13240 0.0003356740
## 
## $calib$probs[[11]]
##         [,1]         [,2]
##   [1,] 11170 0.0002431667
##   [2,] 11175 0.0003058674
##   [3,] 11180 0.0004173059
##   [4,] 11185 0.0004912216
##   [5,] 11190 0.0006095934
##   [6,] 11195 0.0008473031
##   [7,] 11200 0.0013189772
##   [8,] 11205 0.0019574846
##   [9,] 11210 0.0024680836
##  [10,] 11215 0.0030455418
##  [11,] 11220 0.0035792463
##  [12,] 11225 0.0038710174
##  [13,] 11230 0.0044040331
##  [14,] 11235 0.0054210355
##  [15,] 11240 0.0079620663
##  [16,] 11245 0.0115490203
##  [17,] 11250 0.0147571191
##  [18,] 11255 0.0163521858
##  [19,] 11260 0.0178875793
##  [20,] 11265 0.0200005348
##  [21,] 11270 0.0214838960
##  [22,] 11275 0.0216901092
##  [23,] 11280 0.0213666675
##  [24,] 11285 0.0211517421
##  [25,] 11290 0.0211639188
##  [26,] 11295 0.0211639188
##  [27,] 11300 0.0211759431
##  [28,] 11305 0.0211076456
##  [29,] 11310 0.0208558354
##  [30,] 11315 0.0201676844
##  [31,] 11320 0.0189907224
##  [32,] 11325 0.0176393838
##  [33,] 11330 0.0165320925
##  [34,] 11335 0.0157655159
##  [35,] 11340 0.0150209931
##  [36,] 11345 0.0136167806
##  [37,] 11350 0.0117561003
##  [38,] 11355 0.0098527346
##  [39,] 11360 0.0087180971
##  [40,] 11365 0.0082582350
##  [41,] 11370 0.0087390584
##  [42,] 11375 0.0095393382
##  [43,] 11380 0.0086182975
##  [44,] 11385 0.0085892187
##  [45,] 11390 0.0097580240
##  [46,] 11395 0.0116720486
##  [47,] 11400 0.0124795988
##  [48,] 11405 0.0112322257
##  [49,] 11410 0.0092572926
##  [50,] 11415 0.0082210782
##  [51,] 11420 0.0073497634
##  [52,] 11425 0.0059603323
##  [53,] 11430 0.0045949143
##  [54,] 11435 0.0044043719
##  [55,] 11440 0.0055059091
##  [56,] 11445 0.0075905987
##  [57,] 11450 0.0092693051
##  [58,] 11455 0.0091613107
##  [59,] 11460 0.0083033494
##  [60,] 11465 0.0084037676
##  [61,] 11470 0.0093580069
##  [62,] 11475 0.0104881346
##  [63,] 11480 0.0114848180
##  [64,] 11485 0.0123917444
##  [65,] 11490 0.0138705951
##  [66,] 11495 0.0145735423
##  [67,] 11500 0.0141869159
##  [68,] 11505 0.0130289508
##  [69,] 11510 0.0108000973
##  [70,] 11515 0.0076218649
##  [71,] 11520 0.0045949143
##  [72,] 11525 0.0032084797
##  [73,] 11530 0.0032666387
##  [74,] 11535 0.0042876637
##  [75,] 11540 0.0066656760
##  [76,] 11545 0.0090359194
##  [77,] 11550 0.0101112881
##  [78,] 11555 0.0098261392
##  [79,] 11560 0.0086970663
##  [80,] 11565 0.0075183601
##  [81,] 11570 0.0072369985
##  [82,] 11575 0.0075802606
##  [83,] 11580 0.0083021900
##  [84,] 11585 0.0101819953
##  [85,] 11590 0.0130867863
##  [86,] 11595 0.0150735716
##  [87,] 11600 0.0153114814
##  [88,] 11605 0.0137137193
##  [89,] 11610 0.0112624504
##  [90,] 11615 0.0088740427
##  [91,] 11620 0.0069238844
##  [92,] 11625 0.0053786917
##  [93,] 11630 0.0038990933
##  [94,] 11635 0.0027487807
##  [95,] 11640 0.0018340155
##  [96,] 11645 0.0012126721
##  [97,] 11650 0.0009088548
##  [98,] 11655 0.0009579867
##  [99,] 11660 0.0012733650
## [100,] 11665 0.0018279635
## [101,] 11670 0.0024253739
## [102,] 11675 0.0029119387
## [103,] 11680 0.0030531606
## [104,] 11685 0.0027568592
## [105,] 11690 0.0024947332
## [106,] 11695 0.0027658880
## [107,] 11700 0.0031014925
## [108,] 11705 0.0028685976
## [109,] 11710 0.0020814492
## [110,] 11715 0.0012733650
## [111,] 11720 0.0007993143
## [112,] 11725 0.0006345468
## [113,] 11730 0.0006029565
## [114,] 11735 0.0005593572
## [115,] 11740 0.0005730325
## [116,] 11745 0.0007209185
## [117,] 11750 0.0009330359
## [118,] 11755 0.0009101168
## [119,] 11760 0.0006148011
## [120,] 11765 0.0003682801
## [121,] 11770 0.0003078592
## [122,] 11775 0.0002788133
## [123,] 11780 0.0002364413
## [124,] 11785 0.0001869770
## [125,] 11790 0.0001693674
## [126,] 11795 0.0001963664
## [127,] 11800 0.0002593792
## [128,] 11805 0.0002945617
## [129,] 11810 0.0003090989
## [130,] 11815 0.0003004862
## [131,] 11820 0.0003393702
## [132,] 11825 0.0003649080
## [133,] 11830 0.0003393702
## [134,] 11835 0.0002510465
## 
## $calib$probs[[12]]
##         [,1]         [,2]
##   [1,] 20430 0.0003858044
##   [2,] 20440 0.0004389171
##   [3,] 20450 0.0005184455
##   [4,] 20460 0.0006280199
##   [5,] 20470 0.0007611816
##   [6,] 20480 0.0009022009
##   [7,] 20490 0.0010847039
##   [8,] 20500 0.0013357687
##   [9,] 20510 0.0016108617
##  [10,] 20520 0.0019306135
##  [11,] 20530 0.0024127213
##  [12,] 20540 0.0029195164
##  [13,] 20550 0.0034282627
##  [14,] 20560 0.0040993416
##  [15,] 20570 0.0047617309
##  [16,] 20580 0.0054759211
##  [17,] 20590 0.0062114183
##  [18,] 20600 0.0069500317
##  [19,] 20610 0.0073913912
##  [20,] 20620 0.0078639612
##  [21,] 20630 0.0081606228
##  [22,] 20640 0.0083696137
##  [23,] 20650 0.0083772675
##  [24,] 20660 0.0084816210
##  [25,] 20670 0.0083867666
##  [26,] 20680 0.0085870878
##  [27,] 20690 0.0085870878
##  [28,] 20700 0.0086936737
##  [29,] 20710 0.0090088377
##  [30,] 20720 0.0094453831
##  [31,] 20730 0.0097918196
##  [32,] 20740 0.0104990857
##  [33,] 20750 0.0111304292
##  [34,] 20760 0.0119264071
##  [35,] 20770 0.0129079039
##  [36,] 20780 0.0140983167
##  [37,] 20790 0.0154089115
##  [38,] 20800 0.0169169267
##  [39,] 20810 0.0188599462
##  [40,] 20820 0.0209890123
##  [41,] 20830 0.0235278218
##  [42,] 20840 0.0260330244
##  [43,] 20850 0.0287487235
##  [44,] 20860 0.0311847511
##  [45,] 20870 0.0332056013
##  [46,] 20880 0.0345646224
##  [47,] 20890 0.0351508650
##  [48,] 20900 0.0349608117
##  [49,] 20910 0.0340288410
##  [50,] 20920 0.0325199154
##  [51,] 20930 0.0304871789
##  [52,] 20940 0.0281533937
##  [53,] 20950 0.0255426969
##  [54,] 20960 0.0229651525
##  [55,] 20970 0.0206390433
##  [56,] 20980 0.0183610050
##  [57,] 20990 0.0162952622
##  [58,] 21000 0.0147934789
##  [59,] 21010 0.0133980851
##  [60,] 21020 0.0123640394
##  [61,] 21030 0.0113091106
##  [62,] 21040 0.0104440318
##  [63,] 21050 0.0096326081
##  [64,] 21060 0.0088725825
##  [65,] 21070 0.0082571769
##  [66,] 21080 0.0075782798
##  [67,] 21090 0.0070327461
##  [68,] 21100 0.0066019549
##  [69,] 21110 0.0061927951
##  [70,] 21120 0.0058799157
##  [71,] 21130 0.0055848379
##  [72,] 21140 0.0053036702
##  [73,] 21150 0.0050293171
##  [74,] 21160 0.0048363024
##  [75,] 21170 0.0045232488
##  [76,] 21180 0.0042846728
##  [77,] 21190 0.0041177396
##  [78,] 21200 0.0038984793
##  [79,] 21210 0.0037455790
##  [80,] 21220 0.0036541486
##  [81,] 21230 0.0036143368
##  [82,] 21240 0.0034737045
##  [83,] 21250 0.0034378911
##  [84,] 21260 0.0034038220
##  [85,] 21270 0.0033594916
##  [86,] 21280 0.0033157349
##  [87,] 21290 0.0032614315
##  [88,] 21300 0.0032187433
##  [89,] 21310 0.0030939949
##  [90,] 21320 0.0030135271
##  [91,] 21330 0.0028214253
##  [92,] 21340 0.0025847174
##  [93,] 21350 0.0023840209
##  [94,] 21360 0.0021497443
##  [95,] 21370 0.0019122653
##  [96,] 21380 0.0016551621
##  [97,] 21390 0.0013883812
##  [98,] 21400 0.0011720554
##  [99,] 21410 0.0009488750
## [100,] 21420 0.0008016994
## [101,] 21430 0.0006714234
## [102,] 21440 0.0005711339
## [103,] 21450 0.0004713771
## [104,] 21460 0.0003911450
## 
## $calib$probs[[13]]
##        [,1]         [,2]
##  [1,] 19510 0.0003377589
##  [2,] 19520 0.0004365525
##  [3,] 19530 0.0005542492
##  [4,] 19540 0.0007247793
##  [5,] 19550 0.0009921475
##  [6,] 19560 0.0013366044
##  [7,] 19570 0.0018006770
##  [8,] 19580 0.0023962901
##  [9,] 19590 0.0032649269
## [10,] 19600 0.0043316368
## [11,] 19610 0.0055675082
## [12,] 19620 0.0069273063
## [13,] 19630 0.0083448239
## [14,] 19640 0.0093602621
## [15,] 19650 0.0102408035
## [16,] 19660 0.0107803598
## [17,] 19670 0.0110566666
## [18,] 19680 0.0109178900
## [19,] 19690 0.0103752849
## [20,] 19700 0.0095991576
## [21,] 19710 0.0087510823
## [22,] 19720 0.0079641716
## [23,] 19730 0.0072220813
## [24,] 19740 0.0067198626
## [25,] 19750 0.0063337611
## [26,] 19760 0.0061393117
## [27,] 19770 0.0060366862
## [28,] 19780 0.0061126382
## [29,] 19790 0.0064802444
## [30,] 19800 0.0071775365
## [31,] 19810 0.0080374255
## [32,] 19820 0.0092301380
## [33,] 19830 0.0109646153
## [34,] 19840 0.0127656660
## [35,] 19850 0.0152076870
## [36,] 19860 0.0178133738
## [37,] 19870 0.0207455222
## [38,] 19880 0.0234324403
## [39,] 19890 0.0258211808
## [40,] 19900 0.0277158453
## [41,] 19910 0.0290520979
## [42,] 19920 0.0296871819
## [43,] 19930 0.0297911697
## [44,] 19940 0.0295768582
## [45,] 19950 0.0291247278
## [46,] 19960 0.0286375951
## [47,] 19970 0.0280326508
## [48,] 19980 0.0275970153
## [49,] 19990 0.0269981417
## [50,] 20000 0.0264518411
## [51,] 20010 0.0258697090
## [52,] 20020 0.0252555036
## [53,] 20030 0.0244484866
## [54,] 20040 0.0236048884
## [55,] 20050 0.0227321976
## [56,] 20060 0.0217459980
## [57,] 20070 0.0206599487
## [58,] 20080 0.0193825169
## [59,] 20090 0.0178133738
## [60,] 20100 0.0163790888
## [61,] 20110 0.0148209189
## [62,] 20120 0.0135001677
## [63,] 20130 0.0122503169
## [64,] 20140 0.0110780418
## [65,] 20150 0.0098613919
## [66,] 20160 0.0088567196
## [67,] 20170 0.0079254274
## [68,] 20180 0.0070681110
## [69,] 20190 0.0063865520
## [70,] 20200 0.0057637740
## [71,] 20210 0.0051964345
## [72,] 20220 0.0046809853
## [73,] 20230 0.0041507687
## [74,] 20240 0.0037342546
## [75,] 20250 0.0033076670
## [76,] 20260 0.0029943840
## [77,] 20270 0.0026532869
## [78,] 20280 0.0023716869
## [79,] 20290 0.0020522616
## [80,] 20300 0.0018192312
## [81,] 20310 0.0015892699
## [82,] 20320 0.0014101032
## [83,] 20330 0.0012518498
## [84,] 20340 0.0011120757
## [85,] 20350 0.0010032216
## [86,] 20360 0.0008796947
## [87,] 20370 0.0007712389
## [88,] 20380 0.0006963953
## [89,] 20390 0.0006202513
## [90,] 20400 0.0005702238
## [91,] 20410 0.0005091027
## [92,] 20420 0.0004486690
## [93,] 20430 0.0004033594
## [94,] 20440 0.0003470479
## [95,] 20450 0.0003053426
## 
## 
## 
## $rng
## [1]  5905 21460
## 
## $th0
## [1] 5361 5379
## 
## $depths
##   [1]  39  40  41  42  43  44  45  46  47  48  49  50  51  52  53  54  55  56
##  [19]  57  58  59  60  61  62  63  64  65  66  67  68  69  70  71  72  73  74
##  [37]  75  76  77  78  79  80  81  82  83  84  85  86  87  88  89  90  91  92
##  [55]  93  94  95  96  97  98  99 100 101 102 103 104 105 106 107 108 109 110
##  [73] 111 112 113 114 115 116 117 118 119 120 121 122 123 124 125 126 127 128
##  [91] 129 130 131 132 133 134 135 136 137 138 139 140 141 142 143 144 145 146
## [109] 147 148 149 150 151 152 153 154 155 156 157 158 159 160 161 162 163 164
## [127] 165 166 167 168 169 170 171 172 173 174 175 176 177 178 179 180 181 182
## [145] 183 184 185 186 187 188 189 190 191 192 193 194 195 196 197 198 199 200
## [163] 201 202 203 204 205 206 207 208 209 210 211 212 213 214 215 216 217 218
## [181] 219 220 221 222 223 224 225 226 227 228 229 230 231 232 233 234 235 236
## [199] 237 238 239 240 241 242 243 244 245 246 247 248 249 250 251 252 253 254
## [217] 255 256 257 258 259 260 261 262 263 264 265 266
## 
## $elbows
##   [1]  39  41  43  45  47  49  51  53  55  57  59  61  63  65  67  69  71  73
##  [19]  75  77  79  81  83  85  87  89  91  93  95  97  99 101 103 105 107 109
##  [37] 111 113 115 117 119 121 123 125 127 129 131 133 135 137 139 141 143 145
##  [55] 147 149 151 153 155 157 159 161 163 165 167 169 171 173 175 177 179 181
##  [73] 183 185 187 189 191 193 195 197 199 201 203 205 207 209 211 213 215 217
##  [91] 219 221 223 225 227 229 231 233 235 237 239 241 243 245 247 249 251 253
## [109] 255 257 259 261
## 
## $K
## [1] 112
## 
## $cK
## [1] 263
## 
## $slumpfree
##   [1]  39  40  41  42  43  44  45  46  47  48  49  50  51  52  53  54  55  56
##  [19]  57  58  59  60  61  62  63  64  65  66  67  68  69  70  71  72  73  74
##  [37]  75  76  77  78  79  80  81  82  83  84  85  86  87  88  89  90  91  92
##  [55]  93  94  95  96  97  98  99 100 101 102 103 104 105 106 107 108 109 110
##  [73] 111 112 113 114 115 116 117 118 119 120 121 122 123 124 125 126 127 128
##  [91] 129 130 131 132 133 134 135 136 137 138 139 140 141 142 143 144 145 146
## [109] 147 148 149 150 151 152 153 154 155 156 157 158 159 160 161 162 163 164
## [127] 165 166 167 168 169 170 171 172 173 174 175 176 176 176 176 176 177 178
## [145] 179 180 181 182 183 184 185 186 187 188 189 190 191 192 193 194 195 196
## [163] 197 198 199 200 201 202 203 204 205 206 207 208 209 210 211 212 213 214
## [181] 215 216 217 218 219 220 221 222 223 224 225 226 227 228 229 230 231 232
## [199] 233 234 235 236 237 238 239 240 241 242 243 244 245 246 247 248 249 250
## [217] 251 252 253 254 255 256 257 258 259 260 261 262
## 
## $slumphiatus
## [1] 176
## 
## $slumpdets
##          labID   age error depth cc delta.R delta.STD t.a t.b
## 1  D-AMS039542  5345    28    39  3       0         0   3   4
## 2  D-AMS039543  5806    33    44  3       0         0   3   4
## 3  D-AMS039544  7259    32    67  3       0         0   3   4
## 4  D-AMS039545  8657    45   106  3       0         0   3   4
## 5  D-AMS039546  9736    35   150  3       0         0   3   4
## 6  D-AMS039547 10211    63   161  3       0         0   3   4
## 7  D-AMS039548 10936    41   162  3       0         0   3   4
## 8  D-AMS039549 10747    49   162  3       0         0   3   4
## 9  D-AMS039550 10590    39   168  3       0         0   3   4
## 10 D-AMS039551 11117    51   195  3       0         0   3   4
## 11 D-AMS039552  9977    45   219  3       0         0   3   4
## 12 D-AMS039553 17353    84   250  3       0         0   3   4
## 13 D-AMS039554 16558    69   262  3       0         0   3   4
## 
## $prefix
## [1] "/home/mat150604/MEGA/MEGAsync/workflow/TT3/cores_bacon/TT3/TT3_112"
## 
## $bacon.file
## [1] "/home/mat150604/MEGA/MEGAsync/workflow/TT3/cores_bacon/TT3/TT3_112.bacon"
\end{verbatim}

\section{Discussion}\label{discussion}

\subsection{Depositional dynamics and associated hydroclimatic changes
during the Last Glacial Termination (20.5--12.8 cal ka
BP)}\label{depositional-dynamics-and-associated-hydroclimatic-changes-during-the-last-glacial-termination-20.512.8-cal-ka-bp}

The ATTL record highlights the climatic shifts at mid-latitude South
America during the Last Glacial Termination (\emph{T1}; ca. 20.5 - 12.5
cal ka BP), showing shallow littoral lacustrine conditions and
high-energy sediment input prior to the first human settlements in
central Chile (Figs. 5 and 7). Drier conditions prevailed between 20.5
and 18.5 cal ka BP as indicated by fine-grained deposition consistent
with greater meltwater drainage and dust input (47--51), whereas low
δ¹³C\_car and δ¹⁸O\_car values reflect reduced evaporation and increased
freshwater input from snowmelt (52--54) (Fig. 5F-H).

During the final stage of the LGM, the ATTL basin was characterized by
the deposition of fine \emph{L8 facies}, with high silt and clay
contents (\textgreater80\%), very low TC (\textless0.1\%), and
relatively elevated TS percentages (Figs. 2, 3, and 4). The unimodal
grain size distribution (\textasciitilde4 μm mode) and low sorting
suggest gradual, year-round sediment input, either via long-distance
aeolian transport (53) or meltwater flows from glaciated areas (48). As
in other glaciated settings, the fine-grained material deposited during
glacial times was likely related to glacier dynamics, aeolian dust
(47--51), and local clays derived from weathered soils in the catchment
(49).

Stratigraphic correlations show that the \emph{L8 facies} can be
correlated with Unit 5 (10) and Member 5 (45, 55) from previously
described Tagua Tagua profiles, dated to ca. 21.5 ± 0.65 and 20.4 ± 1.3
cal ka BP, respectively (Fig. 1c). Researchers originally interpreted
these deposits as a massive clay layer rich in \emph{Nothofagus
dombeyi/obliqua} type pollen, suggesting cold and wet conditions or
shallow lacustrine environments. However, the limited presence of
planktonic diatoms, aquatic pollen, and phytoliths, along with only
minor benthic diatoms and freshwater algae, points instead to a shallow
littoral lake setting under persistent, fine detrital input influenced
by a glaciated watershed.

A brief high-energy alluvial event at \textasciitilde19.5 cal ka BP did
not significantly elevate lake levels or bioproductivity (Figs. 3 and
5F-H). Varela (45) describes subvertical cracks filled with authigenic
magnesium oxide and iron minerals in Member 5, consistent with those in
the \emph{L8} and \emph{L7 facies} (Unit 8). Subvertical cracks filled
with authigenic minerals can be interpreted as desiccation cracks,
sub-aquatic sedimentary fractures, or fragipan features (56--59). These
cracks, lacking mud curls and filled with magnesium oxide, iron, and
sulfur minerals, suggest salinity fluctuations or anoxic conditions
promoted shrinkage through deflocculation or intracrystalline dewatering
(56,57).

The high sulfur content of these facies suggests methane emissions and
anaerobic sulfate or metal reduction by methanotrophic archaea (57),
reinforcing the notion of a littoral depositional environment strongly
shaped by watershed inputs.

Our interpretation contrasts with earlier suggestions of cold and wet
conditions at the end of the LGM (10, 44, 45, 65). \emph{Heusser} (44)
proposed the presence of a closed-canopy \emph{Nothofagus} forest
between 28 and 11 ka BP, suggesting cool and wet conditions potentially
driven by a northward shift of the SWW. Recalibrated pollen
chronologies, however, suggest that the percentages of the
\emph{Nothofagus antarctica/dombeyi} type declined earlier than
previously reported (Fig. S6), occurring between approximately 19--16
cal ka BP (initially dated at 14.5 ± 0.35 cal ka BP; Unit 2a, Figs. 7h
and S6) in \emph{Heusser} (1983) and between 14.9--12.7 cal ka BP
(initially dated at 11.71 ± 0.43 cal ka BP; Unit 3, Fig. 1) in
\emph{Valero-Garcés et al.} (2005).

In the central Chile valley, \emph{Nothofagus dombeyi/obliqua} exhibits
a more restricted distribution compared to other members of the genus
and subgenus \emph{Lophozonia} (66). Germination studies reveal that
these species are not strictly dependent on high moisture levels,
challenging the assumption that their expansion was moisture-driven (67,
68). \emph{León-Lobos} and \emph{Ellis} (2007) found that the
germination of \emph{Nothofagus} species (\emph{N. alpina}, \emph{N.
dombeyi}, \emph{N. leonii}, \emph{N. glauca}, and \emph{N. obliqua})
varies more significantly with temperature, suggesting that moisture was
important but not critical for their expansion into central Chile.

Intensified alluvial activity, associated with more frequent and severe
storms at ATTL, is practically coeval with the HS1 (37, 69), CAPE I (23,
70--72), and the ACR/Bölling--Alleröd (36) (Fig. 7), during a period
characterized by increasing winter insolation and decreasing summer
insolation (Figs. 7D, 7E, 7O, and 7P). These changes have been linked to
shifts in global temperature gradients driven by residual Milankovitch
effects (precession and obliquity; Figs. 7A and 7B). Some authors have
posited that these shifts may have triggered a southward migration of
the ITCZ and SWW, a weakening of the American Monsoon System, and
modifications to the PSJ (33, 73--75).

This atmospheric configuration, coupled with the reactivation of
ENSO-like conditions ca. 16 cal ka BP (Figs. 7J and 7Q), likely
amplified the frequency and intensity of alluvial events, displaying
atmospheric patterns similar to those observed in simulations for the
LGM (Fig. 6G). Marine sediment cores off the Chilean coast date
\emph{T1} to between 19.2 and 18.5 cal ka BP, indicating arid conditions
(73). Similarly, Andean glaciers (33°S - 36°S) retreated between 28.0
and 18.0 ka BP (Fig. 7I) without subsequent re-advances during \emph{T1}
(11, 76, 77). Changes in leaf wax δD and δ¹³C isotopes from marine
sediment cores suggest hydrological changes influenced by precessional
cycles, with SWW retreating during \emph{T1} (Figs. 7N and 7G), reaching
a minimum at 11 cal ka BP and again between 5--6 cal ka BP (33). Marine
records also show a 6°C--7°C rise in SST between 19 and 12.5 cal ka BP,
interrupted by cooling around 14.8--13.3 cal ka BP, corresponding to the
ACR (33, 73, 78, 79).

The decline in high-energy processes after ca. 14.8 cal ka BP was
accompanied by fluvial and fine-grained sediment deposition (Fig.
7J-K-L). Orbital changes alone cannot fully explain these phenomena.
Instead, multiple high-magnitude events associated with rapid coupled
oceanic--atmospheric changes likely interacted with gradual orbital
processes to produce the alluvial and fluvial deposits in TT19-3A
between \textasciitilde18.5 and 14.5 cal ka BP. Further studies are
required to better understand the interplay among these factors and
their implications for regional paleoecology.

\section{Changing lake environments, emerging habitats, and first human
settlement (12.9--8.6 cal ka
BP)}\label{changing-lake-environments-emerging-habitats-and-first-human-settlement-12.98.6-cal-ka-bp}

Between ca. 12.8 and 12.4 cal ka BP (\emph{Unit 5}), coarser \emph{L5b}
and \emph{L5a facies} (with a grain size mode at \textasciitilde700 µm;
\emph{rEM5}) indicate high-energy conditions and stronger alluvial
influence in the lake shoreline (Figs. 2, 4, and 5). The increasing
ubiquity of the diatom species \emph{Anomoeoneis sculpta} (Figs. 3 and
S4) ---known for inhabiting alkaline (1350 mg/L CaCO₃), shallow
freshwater environments with elevated pH (9.3) and conductivity (5760
µmS/cm) (80, 81)--- and the rise of planktonic diatoms (e.g., small
fragilarioids) suggest a permanent, alkaline lake established in the
basin (Fig. S4). The presence of \emph{Sporormiella} spores,
megaherbivore remains, and evidence of early human occupation around
12.5 cal ka BP coincides with this transitory alkaline lake phase (22).

From 12.5 to 11.2 cal ka BP (\emph{Unit 5}), a significant depositional
and biogeochemical shift occurred. Greenish to gray sandy silts with
coarser GSDs (\emph{rEM2; L4b facies}) had abundant gastropods and
charcoal spicules, higher carbonate (TIC) content, more positive
δ¹³C\_car and δ¹⁸O\_car, increased diatom abundance (especially \emph{A.
sculpta}), and a rise in woody, shrub, and \emph{Jubaea chilensis}
phytoliths (Figs. 2, 3, 4, and 5). This suggests a moderately alkaline,
shallow lake environment characterized by higher bioproductivity, more
stable water levels, and warmer, more humid conditions. In general,
freshwater lakes with high conductivity and salinity levels exhibit a
negative correlation with diatom diversity (81). The presence of
\emph{J. chilensis} corresponds to the endemic Chilean palm, currently
distributed between \textasciitilde30°S and 35°S (82). Their occurrence
near the 1.65 m depth (ca. 12.4 cal ka BP; see Fig. 4) may indicate
increased humidity, reduced megafaunal pressure, or human transport for
food, temporary habitation, or other uses (83, 84). Alternatively, other
taphonomic and sedimentary transport processes associated with a rising
lake level could also explain its presence.

A brief interval between \textasciitilde11.4 and 11.2 cal ka BP
(\emph{L4a facies}) shows even higher TIC, δ¹³C\_car, δ¹⁸O\_car, and
\emph{rEM2} values, as well as a shift in diatom assemblages toward
tychoplanktonic small fragilarioids and \emph{Aulacoseira granulata}.
These changes suggest a decrease in salinity and an increase in lake
turbulence (Figs. 2, 3, and 4). During this period, freshwater algae
dominate the palynomorph assemblage, and woody and shrub taxa appear
alongside \emph{Poaceae}, \emph{Bambusoideae}, and \emph{Chloridoid}
phytoliths (Fig. 3).

Isotope signatures support warmer conditions and higher lake levels at
the onset of the Holocene. In ATTL sediments, calcite is the principal
carbonate phase. Elevated δ¹³C\_car (\textasciitilde1‰) and δ¹⁸O\_car
(\textasciitilde-0.2‰) values in authigenic calcite indicate
evaporation-driven isotopic enrichment of dissolved inorganic carbon
(DIC). This enrichment, facilitated by isotopic equilibrium exchange
with CO₂, is regulated by temperature, the isotopic composition of
inflows, pH, and lacustrine bioproductivity (10, 85--87). Currently,
there is no data to validate the isotopic enrichment model for meteoric
precipitation proposed by \emph{Valero-Garcés et al.} (2005) for ATTL.
However, \emph{McCormack et al.} (2019) suggest that δ¹⁸O\_car and
δ¹³C\_car values are closer to freshwater runoff compared to other
carbonates like aragonite or dolomite. The isoscape from Chilean lakes
serves as an analogue for time-averaged local precipitation isotopes
where direct data are lacking (88). In central Chile, meteoric water
δ¹⁸O averages from -4‰ to -2‰ in summer and -7.6‰ to -5.9‰ in winter
(88, 89). Data from the \emph{HOLOCHILL} project show notable isotopic
variability in lakes by altitude (Fig. 10S). Coastal areas, such as Lago
Vichuquén, exhibit δ¹⁸O values ranging from -2‰ to 0‰, while Laguna
Matanzas shows values between 7‰ to 8‰. In contrast, mountain lakes like
Laguna del Maule and Laguna Lo Encañado have δ¹⁸O values of
\textasciitilde-11‰ to -12‰.

Altogether, these data suggest a rise in lake levels under warm and
humid conditions at the onset of the Holocene, consistent with an
environment transitioning toward a new ecological state (22). From 11.2
to 8.6 cal ka BP (\emph{Unit 4}), a series of fluctuations in water
levels, detrital input, and carbonate isotopic signals are documented by
several proxies (Figs. 2 and 3).

\emph{L3a-d facies} show alternating clayey silty sands and sandy silts
with variable TIC, δ¹³C\_car and δ¹⁸O\_car values, as well as shifts in
diatom communities and phytolith assemblages (Fig. 3). \emph{L3d facies}
(ca. 11.2 and 10.4 cal ka BP) show a shift in the isotopic signal of
authigenic calcite, accompanied by an abrupt increase in \emph{rEM5}
percentages. Episodes of high-energy alluvial transport, variations in
the presence of \emph{A. sculpta} and small fragilarioids, and frequent
planktonic diatoms reflect changing littoral environments and variable
lake levels (Figs. 3 and S4). Small fragilarioids ---facultative
planktonic organisms from lentic and deltaic environments--- and
\emph{A. granulata} indicate a rise in lake levels, reduced light
availability, and increased wind or tidal turbulence, consistent with
high \emph{rEM5} values (90, 91). Woody and shrub phytoliths, along with
\emph{Poaceae} and \emph{Bambusoideae}, reflect nearby vegetation
responses to these hydrological and climatic oscillations (Figs. 3).

The \emph{L3c facies} (10.4 and 9.8 cal ka BP) show variable percentages
of \emph{rEM2} and \emph{rEM3}, accompanied by elevated \emph{rEM4}
values, which highlight high variability of alternating sedimentary
processes (littoral, fluvial, aeolian). During deposition of \emph{L3b
facies} (between ca. 9.8 and 9.0 cal ka BP), changes in the fluvial
sediment delivery (\emph{rEM3}) suggest some hydroclimatic instability.
\emph{A. sculpta} increases rapidly in this facies, while small
fragilarioids practically disappear, indicating a shallow to ephemeral
and subaerial depositional environment.

The phytolith assemblage remains virtually unchanged. Finally,
deposition of \emph{L3a facies} (9.0 - 8.6 cal ka BP), characterized by
a high content of fine sand (\emph{rEM4}), and a decrease in planktonic
diatom species, culminated in the Early Holocene. Stable lake levels,
finer sediments (\emph{rEM2} and a minority of \emph{rEM1}), a dominance
of hygro-hydrophilous communities, and elevated organic productivity
(\emph{Units 3} and \emph{2}, ca. 8.6 - 6.0 cal ka BP) fostered greater
diatom diversity and more mesotrophic, alkaline conditions. Finer
sedimentation and sustained presence of planktonic diatoms
(\emph{Stephanodiscus meneghinianus}) suggest relatively higher lake
levels compared to \emph{Unit 4}, and less sediment input from the
watershed. At the time of deposition of \emph{Unit 3} (8.6 - 6.5 cal ka
BP), the later occupation of the archaeological site TT-3 occurred.

\section{Early human-ecosystem interactions in ATTL during the late
Pleistocene/Holocene
transition}\label{early-human-ecosystem-interactions-in-attl-during-the-late-pleistoceneholocene-transition}

The progressive environmental transformation toward the end of the
Pleistocene and the onset of the Holocene ---evidenced by rising lake
levels and the diversification of available resources (22)---
represented a turning point for human occupation in the ATTL. These
changes broadened subsistence possibilities for the earliest
hunter-gatherers, marking the onset of a distinct ecological stage in
South America (Fig. 7F).

In the ATTL, the transition from the Pleistocene to the Holocene (14.5
to 11.7 cal ka BP) is characterized by a gradual rise in lake levels
(Figs. 5D and 7L). An alluvial event between ca. 12.8 and 12.5 cal ka BP
and the development of beach-like depositional environments in littoral
areas indicate a return to drier conditions (Figs. 5 and 7). The end of
this event coincided with the discovery of a combustion feature along
with megafaunal remains, dated to 12.5--12.4 cal ka BP (22).

Archaeological evidence indicates that human occupation of the ATTL
began during the Fishtail point horizon (92). The archaeological sites
in the ATTL (\emph{TT-1}, \emph{TT-2}, and \emph{TT-3}), dated to ca.
12.6--11.6 cal ka BP, provide unequivocal evidence of direct
human--megafauna interactions, including gomphotheres, horses, and
cervids (22, 93--96). These sites also document the consistent
exploitation of small game (primarily birds and frogs) and the use of
local plant resources (22, 93--96). Interpreted as logistic or ephemeral
camps, these settlements were part of a broader and more complex
mobility system that included at least the Andean range (22, 94).
Regionally, however, late Pleistocene settlements are conspicuously
absent, whereas Early Holocene archaeological sites are common in the
Andes, along the coastal range, and the coastline (42).

Amid these regional cultural developments, paleoclimatic data suggest
increasing water availability and temperatures, coupled with a reduced
influence of the SWW near 30°S at the end of the Late Pleistocene (Figs.
2, 3, 7L-M-N). This was followed by a northward shift of the SWW at the
beginning of the Holocene, along with large-amplitude REC associated
with YD and ENSO-like phenomena (33, 40, 97). These
processes---partially driven by obliquity-induced orbital
changes---likely governed the variability of both the SWW and the SPSA
in central Chile between 12.5 and 11 cal ka BP, forming a complex
interplay not fully captured by other records at this latitude (12, 33,
73, 75).

Against this dynamic environmental backdrop, the earliest human groups
present at ATTL adopted logistical settlement patterns under
increasingly favorable conditions, taking advantage of a wide range of
lacustrine resources. This strategy persisted until around 11--9 cal ka
BP (22, 95). The intertwined development of climate, megafauna, and
human strategies underscores the profound impact that shifting
environmental regimes had on the cultural trajectories of early
hunter-gatherer communities.

\section{Early-Middle Holocene: increased environmental variability and
emerging human settlement
patterns}\label{early-middle-holocene-increased-environmental-variability-and-emerging-human-settlement-patterns}

Early and Middle Holocene environmental reconstructions in central Chile
suggest a dynamic interplay among regional teleconnections,
climate-driven erosive processes, tectonics, volcanism, and the evolving
impacts of early hunter-gatherers on the ecosystem (10, 44, 98--101).
The TT19-3A record indicates higher lake levels, enhanced evaporative
regimes, and the emergence of diverse ecological assemblages during the
Early Holocene compared to the Late Pleistocene (Figs. 5 and 7). The
higher abundance of woody dicotyledon phytoliths indicates the
prevalence of sclerophyll forests, whereas the occurrence of
\emph{Jubaea chilensis} highlights the abundance of palm trees, likely
more prevalent due to the absence of megaherbivores, humid climatic
conditions, and/or human use (84). Phytoliths from \emph{Poaceae} and
\emph{Pooideae} groups indicate the coexistence of grassland and
sclerophyll shrubs. An increase in \emph{Bambusoideae} morphotypes
suggests the expansion of \emph{Chusquea} sp., while the presence of
\emph{Chloridoid} phytoliths---commonly found in arid or semi-arid
regions---points to possible drought periods.

Between 11.5 and 11 cal ka BP, the proportion of fluvial deposits
reached its highest levels (Figs. 5D and 7L), possibly triggered by a
northward shift of the SWW (33) and weakening of the SPSA (Figs. 7N and
7G). After this peak fluvial input, the record indicates shifts in
high-energy sedimentary transport and an increase in diatoms (small
fragilarioids) associated with turbulent lacustrine environments between
11 and 8.5 cal ka BP (Figs. 7J-L). It is possible that high ENSO-like
activity along the Peruvian coast (Fig. 7R) aligns with the changes
recorded in the ATTL and drought periods at Laguna Aculeo (98, 99). Such
large-scale events might have triggered a scaling break effect,
potentially strengthening STJ conditions and contributing to drought
events (19, 102). Variability in STJ intensity, influenced by ENSO-like
phenomena and possibly combined with residual effects of
glacial-interglacial orbital changes (Milankovitch cycles), may have
contributed to the formation of alluvial and fluvial deposits observed
during this period in central Chile (15, 17, 103).

During the first half of the Mid Holocene (8.5--6 cal ka BP), central
Chile experienced a gradual shift toward increased aridity---though less
extreme than events recorded between 11 and 8.5 cal ka BP (Figs. 5,
7H-J). This is supported by PMIP4-CMIP6 multi-model ensemble outputs
(Figs. 6, 11S-13S; Table S2), which indicate a slight increase in zonal
winds in central Chile under wet conditions compared to pre-industrial
levels (Fig. 6F). At lower latitudes, negative zonal wind anomalies
suggest La Niña-like or ENSO-neutral conditions, linked to drier periods
in central Chile (13, 103). Wet conditions, however, show a slight
increase in precipitation relative to the pre-industrial period, while
dry conditions exhibit precipitation and zonal wind patterns similar to
pre-industrial levels (Figs. 6D and 6H). Although no evidence suggests a
northward displacement of the SWW during the Mid Holocene, the TT19-3A
record indicates higher fluvial inputs and elevated lake levels compared
to other central Chilean lacustrine records from the same period (98,
104).

The Early Holocene begins with an important change in the
hunter-gatherer settlement pattern at the ATTL but also in the
surrounding areas. Late Pleistocene logistical and/or ephemeral bases
are replaced by more intensive residential camps at ca. 9.0 cal ka BP
(105, 106). This shift coincided with the second phase of human
occupation identified at \emph{TT-3} (22) in the lower portion of the
\emph{L2 facies} (ca. 8.5 cal ka BP). Nevertheless, considering the
available radiocarbon dates, the Mid Holocene (between ca. 7.2--6.0 ka
cal BP) was possibly the most intensive period of human occupation of
the ATTL (22, 105).

Human burials were added into residential areas, forming complex mounds
characterized by intensive domestic refuse and mixed stratigraphic
deposits. Some isolated dates suggest that these sites were probably
occupied less recurrently until the beginning of the Late Holocene (ca.
4.0 cal ka BP). At least seven of these structures were documented at
the lake's shore, but only a few were systematically excavated. Thus,
their synchronicity, anthropological, artefactual, and ecofactual
particularities, and their possible relations are not yet fully
understood. The systematic presence of human burials in these sites
stresses the pivotal role that ATTL played for the hunter-gatherers
during climatic oscillations, and allows us to hypothesize mobility
reduction and, possibly, incipient territorialism.

On a regional scale, human concentration in particularly productive
locations such as the ATTL could explain the scarce record of Mid
Holocene archaeological sites elsewhere in central Chile. In this
scenario, the lake acts as a refuge by concentrating abundant resources
such as water, food, and raw materials, especially during periods of
prolonged drought (106, 107).

\section{Implications and
Limitations}\label{implications-and-limitations}

The ATTL basin in central Chile offers a detailed record of climatic and
human dynamics during the Late Pleistocene and early Holocene. The
integration of stratigraphic, geochemical, and paleoecological data
highlights the role of local hydroclimatic changes in guiding early
human mobility, resource use, and settlement choices. As lake levels
fluctuated and resource availability shifted, hunters and gatherers
appear to have adapted by altering subsistence strategies and occupation
intensity.

Coupling local lacustrine evidence with large-scale climate models
(PMIP4-CMIP6) reveals that regional drivers (SWW, SPSA, STJ, ENSO-like
phenomena) manifested in highly localized ways. This underscores the
importance of linking local data sets to global circulation patterns to
avoid oversimplification of paleoclimate scenarios. During the Late
Pleistocene, high-energy fluvial inflows driven by atmospheric shifts
(STJ intensity, ITCZ, and SWW displacement) occurred, but these changes
did not significantly impact the lake's water balance, indicating a
complex hydrological regime needing further study. These shifts coincide
with global climate phenomena such as the HS1, B/A, and the CAPE and ACR
events in South America.

The Pleistocene-Holocene transition brought warmer, wetter conditions,
driving rapid environmental changes and the development of new habitats
in the lake and its watershed. The clear association of human presence
with megafaunal remains and Sporormiella spores indicates that shifting
shoreline habitats, prey distribution, and hunter-gatherer activities
converged at ATTL. This study provides valuable insights into how
climate and humans jointly influenced late Pleistocene megafauna
extinctions and dispersals. Rapid lake-level fluctuations and episodes
of high-energy deposition during the Pleistocene and Holocene serve as
analogs for modern climate variability. Historical ecology, including
the use of refugia and settlement reorganization, offers valuable
insights into resilience strategies for current and future water
resource planning.

Evidence from the ATTL basin indicates that, after a humid Early
Holocene phase, increased depositional variability and shoreline habitat
changes were driven by ENSO-like shifts. However, during the Mid
Holocene, the ATTL shows relatively stable lake levels---contrasting
with other central Chilean sites---thereby providing favorable
conditions for humans and fauna. This hydrological stability, coupled
with abundant resources, suggests that the ATTL basin functioned as an
ecological refuge during drought periods. The intensification of
occupation, formation of mounds, and emergence of more complex social
practices (e.g., burials) reflect how favorable local conditions could
catalyze socio-cultural development.

The TT19-3A sequence underscores the strong interplay between climate
and human history in central Chile, showing how early hunter-gatherers
adapted to periods of both abundance and stress. While the
``ecorefugio'' model derived from this study provides valuable insights,
caution is warranted when applying it to other South American contexts,
where variations in geology, topography, and climate may yield distinct
sedimentological and ecological outcomes. Although robust age-depth
models have been used, certain stratigraphic units---particularly those
from the late Pleistocene---remain less densely dated, potentially
obscuring abrupt environmental fluctuations and complicating the
correlation of short-lived events with archaeological contexts.

Further analyses---such as SEM-based quartz grain microtexture studies,
multi-element isotopes, or modern data (e.g., clumped isotope
(Δ\textsuperscript{47}) carbonate, aDNA, or GDGTs/δD in organic
biomarkers)---would strengthen hydroclimate interpretations and help
validate assumed isotope enrichments (δ\textsuperscript{13}C and
δ\textsuperscript{18}O). Additionally, global climate models, while
providing important large-scale insights, often lack sufficient temporal
resolution and fail to capture localized phenomena (e.g., convective
storms, orographic precipitation) at mid-latitudes, making it difficult
to align their outputs with high-resolution sedimentary records.

Finally, lake-level changes, erosional reworking, and post-depositional
processes can skew the archaeological record in shoreline contexts,
causing certain intervals (e.g., the early Holocene) to appear less
populated due to sedimentary discontinuities rather than true
occupational gaps. Broader field surveys and correlating multiple cores
would help confirm and refine interpretations of human settlement
intensity throughout the region.

A high-energy depositional phase occurred between 18.5 and 16.5 cal ka
BP, marked by sandy \emph{L7 facies} (about 70\% - 80\% sand) with low
TC values, and the lowest δ¹³C\_car and δ¹⁸O\_car values in the record
(Figs. 2 and 4). The coarse, bimodal grain size distribution
(\emph{rEM4}; main peak \textasciitilde420 µm, secondary
\textasciitilde60 µm, see Fig. S3) suggests heavy rainfall/intense
storms, or rapid melt-driven discharge that delivered coarse sediments
into the lake's littoral zone (45, 47, 59). The lowest δ¹⁸O\_car values
likely reflect colder waters, an increase in the amount effect, or a
different moisture source (51), whereas lower δ¹³C\_car values possibly
reflect increased CO₂ input from terrestrial respiration or lower
aquatic productivity.

Later, from 16.5 to 14.5 cal ka BP, a shift to coarser sediments and
stormier conditions occurred again without increasing lake levels (Figs.
5a, b, e). This interval corresponds to a period of intense hydrological
and depositional variability, linking local sedimentary patterns to the
broader regional climate context.

To understand the atmospheric patterns during this period, we analyzed
precipitation and zonal wind outputs from the PMIP4-CMIP6 multimodel
ensemble simulations for the LGM (Figs. 6, 7S--9S; Table S2). The
multimodel ensemble indicates that during the LGM, dry conditions were
associated with intensified zonal wind speeds at low and mid-latitudes,
accompanied by negative anomalies at subtropical and high latitudes
(Figs. 6a, 6c). In contrast, wet anomalies are linked to weakened
tropical trade winds, akin to warm ENSO-like conditions, accompanied by
positive anomalies along extratropical storm tracks (Figs. 6E, 6G).
Zonal wind speeds at tropical and subtropical latitudes show positive
anomalies over the southern Pacific, while negative anomalies prevail
over the central-southern Andes. A similar atmospheric pattern to those
of the glacial period likely contributed to extreme events between 18.5
and 14.5 cal ka BP (Figs. 5a, 5b). Sedimentological and limnological
evidence, however, suggests this was insufficient to maintain elevated
lake levels.

Between 16.5 and 12.8 cal ka BP, sedimentation reverted to sandy silts
(Unit 7, \emph{L6a} and \emph{L6b facies}), with varying clay
percentages and low TS and TC values. Microfossils remained with low
abundance, except for two samples showing a temporary increase in
\emph{Poaceae} phytoliths, benthic diatoms, and \emph{Cyperaceae},
hinting at short-lived stabilization of littoral and riparian
communities (Figs. 3 and S4). The presence of freshwater algae
(\emph{Botryococcus} and \emph{Spirogyra}), littoral vegetation such as
\emph{Typha}, and sporadic fern spores (\emph{Pteris}) indicates
wetlands and riparian areas adjacent to the basin (Fig. S5). δ¹³C\_car
and δ¹⁸O\_car values rose relative to Unit 8, indicating changes in
hydrological balance (precipitation - evaporation balance), seasonality,
or source area.

From 16.5 to 14.5 cal ka BP (\emph{L6b facies}), grain size
distributions (\emph{rEM5}) were bimodal (\textasciitilde700 μm and
\textasciitilde65 μm), more consistent with fluvial or reworked
fluvial/lacustrine-aeolian processes (47--51, 59). Such distributions
occur where braided rivers and fluvial processes overlap, mixing coarse
and fine sediments (49, 60). Proximal aeolian coarse-grained silt
transport typically shows a unimodal grain-size distribution between 40
and 90 μm (49, 60). In contrast, GSDs exceeding a \textasciitilde250 μm
mode suggest high-energy fluvial processes. Moreover, in other
environments and depositional process simulations, grain sizes ranging
from 250 to 850 μm have been linked to intense runoff, alluvial fan
formation, and are also observed in sandy coastal settings as well as
wind-eroded sediments (61--64).

After 14.5 cal ka BP (\emph{L6a facies}), a shift toward finer particles
(\emph{rEM1}, \textasciitilde3.8 µm; \emph{rEM2}, \textasciitilde60 µm)
suggests improved hydrodynamic sorting, potentially linked to rising
lake levels or greater availability of finer sediments. These conditions
reflect a dynamic interplay between wind-driven, glacial-mediated, and
fluvial processes influencing deposition at the ATTL site during glacial
time. Further analysis of the quartz surface texture is necessary to
differentiate between aeolian and fluvial processes.

\section{Materials and Methods}\label{materials-and-methods}

In November 2019, we recovered a sedimentary sequence measuring 15 cm in
length, 10 cm in width, and 280 cm in depth from the exposed profile of
the eastern wall of a 1×2m excavation, corresponding to \emph{Unit 1} of
the geoarchaeological test pits at TT-3 site (\emph{TT19-3A-1E}, Fig.
1). Sediment samples, between 50 and 100g of material, were obtained at
1-cm intervals along the sequence (see Supplementary Methods for more
details). The uppermost 8 samples (0 - 0.38 m) were discarded due to
sediment alteration caused by modern agricultural practices.

Sixty-six samples taken every 2cm were analyzed for total carbon
(\emph{TC}), total inorganic carbon (\emph{TIC}), total organic carbon
(\emph{TOC=TC−TIC}), total nitrogen (\emph{TN}), and total sulfur
(\emph{TS}) using a LECO SCN furnace at the Pyrenean Institute of
Ecology (\emph{IPE-CSIC}) in Spain. Stable isotope compositions
(δ¹³C\_car and δ¹⁸O\_car) of the carbonate fraction from these samples
were determined using a Gas Bench system interfaced with a Delta XP mass
spectrometer (Thermo Finnigan, Bremen, Germany) at the Andalusian
Institute of Earth Science (\emph{CSIC}) in Granada, Spain.

The GSDs (from 0.04µm to 2000µm) were measured in triplicate using a
Malvern Mastersizer 2000 APA2000 with Hydro 2000G AWA2000 at
\emph{IPE-CSIC}. Bulk sediment samples were pre-treated to remove
organic matter and carbonates. Forty sediment samples were prepared for
palynomorph extraction following standard chemical procedures (65, 108).
Due to taphonomic and preservation issues, terrestrial pollen is highly
underrepresented; therefore, only pollen from hygro-hydrophytes, fern
spores, and other NPPs have been analyzed (108--110).

The extraction of phytoliths and diatoms was conducted on twenty
sediment samples, which were processed at the Archaeobotany and
Palaeoecology Laboratory of the University of Exeter. A protocol based
on standard procedures for phytolith extraction was employed, using
approximately 5 g of sediment. We removed carbonates using 36\% HCl and
digested organic materials with 70\% HNO₃. Phytoliths and diatoms were
separated through flotation using ZnBr₂ prepared to a density of 2.3
g/cm³. We counted phytoliths, analyzing up to 200 specimens per sample,
or the total number of specimens examined in cases where the slide
contained fewer than 200. Finally, the phytoliths were classified
according to the International Code for Phytolith Nomenclature
(111--113).

Diatom valves on Naphrax-mounted slides were observed using a Zeiss
Axioscope A1 stereomicroscope with an oil immersion objective of 1000x.
A minimum of 400 valves were counted when possible, identifying to
species level with the exception of small fragilarioids that require
higher magnification (114).

We reconstructed the TT19-3A chronology using fourteenth AMS radiocarbon
dates on charcoal and bulk sediments (Table S1).

\subsection{Statistical Analysis}\label{statistical-analysis}

Radiocarbon ages were calibrated using the Southern Hemisphere
calibration curve (115), and age-at-depth assemblage was established
using the \emph{Bacon} and \emph{Bchron} models implemented in the
\texttt{geochronR} package (116--118). Robust End-member Mixing Analysis
(\emph{rEMMA}) was performed using the R package \texttt{EMMAgeo} v0.9.8
to infer sediment sources, transport, and depositional processes (117,
118).

Sedimentary facies were characterized based on macroscopic and
microscopic descriptions (120), and artificial facies clusters were
identified using machine learning techniques (121). The analysis
employed an unsupervised classification algorithm, utilizing geochemical
data (\emph{TOC}, \emph{TIC}, \emph{TN}, \emph{TS}, \emph{GSDs}, and
stable isotopes) implemented in the \texttt{scikit-learn} library (122)
in Python v3.9. Stratigraphic diagrams and timeseries were analyzed and
plotted using \texttt{tidypaleo} v0.1.1, \texttt{actR} v0.2.0,
\texttt{ggplot} v3.3.5, and \texttt{rioja} v0.9-26 packages in R 4.2.2
software.

Circulation anomalies for the Mid Holocene and LGM periods were analyzed
using PMIP4 ensemble simulations from the \emph{MPI-ESM1-2-LR},
\emph{MIROC-ES2L}, \emph{CESM2-WACCM-FV2}, \emph{AWI-ESM1-2-LR}, and
\emph{CESM2} models. These analyses focused on global changes in zonal
wind at 500 hPa and total precipitation relative to the pre-industrial
period (\emph{piControl}). To define wet and dry events in central Chile
(75°W--70°W, 32°S--40°S), a Regional Precipitation Index (\emph{RPI})
was applied, defining wet events as \emph{RPI\textgreater120\%} and dry
events as \emph{RPI\textless80\%} (see Supplementary Methods).

Using the \texttt{climlab} and \texttt{pyleoclim} libraries in Python
v3.9, we analyzed past orbital parameters and paleoclimate data through
wavelet and spectral analysis (Fig. 7Q-R) (123). Human activities were
assessed at the local scale through the analysis of lithic components
found at TT-3, and at the regional scale using the \emph{p3k14c}
database (46) of archaeological radiocarbon dates.

We calculated summed probability distributions (\emph{SPD}) of
archaeological radiocarbon dates using the \texttt{rcarbon} v1.5.1
package in R 4.2.2 for South America and central Chile (between 32° and
37°S). Environmental and orbital changes from paleoclimate records and
models were summarized using available literature.

\section{References}\label{references}

\begin{enumerate}
\def\labelenumi{\arabic{enumi}.}
\tightlist
\item
  Barnosky, A. D., Koch, P. L., Feranec, R. S., Wing, S. L. \& Shabel,
  A. B. Assessing the Causes of Late Pleistocene Extinctions on the
  Continents. Science 306, 70--75 (2004).
\item
  Mann, D. H., Groves, P., Gaglioti, B. V. \& Shapiro, B. A.
  Climate-driven ecological stability as a globally shared cause of Late
  Quaternary megafaunal extinctions: the Plaids and Stripes Hypothesis.
  Biol. Rev.~94, 328--352 (2019).
\item
  Gutiérrez, E. G. \& Stockli, D. F. Fluvial response to Late
  Pleistocene--Holocene climate change in the Colorado River drainage,
  central Texas, USA. Geology 51, 929--934 (2023).
\item
  Braje, T. J. \& Erlandson, J. M. Human acceleration of animal and
  plant extinctions: A Late Pleistocene, Holocene, and Anthropocene
  continuum. Anthropocene 4, 14--23 (2013).
\item
  Mottl, O. et al.~Global acceleration in rates of vegetation change
  over the past 18,000 years. Science 372, 860--864 (2021).
\item
  Andermann, T., Faurby, S., Turvey, S. T., Antonelli, A. \& Silvestro,
  D. The past and future human impact on mammalian diversity. Sci. Adv.
  6, eabb2313 (2020).
\item
  Malhi, Y. et al.~Megafauna and ecosystem function from the Pleistocene
  to the Anthropocene. Proc. Natl. Acad. Sci. 113, 838--846 (2016).
\item
  Baumann, C. The paleo-synanthropic niche: a first attempt to define
  animal's adaptation to a human-made micro-environment in the Late
  Pleistocene. Archaeol. Anthropol. Sci. 15, 63 (2023).
\item
  Mayle, F. E. \& Iriarte, J. Integrated palaeoecology and archaeology
  -- a powerful approach for understanding pre-Columbian Amazonia. J.
  Archaeol. Sci. 51, 54--64 (2014).
\item
  Valero-Garcés, B. L. et al.~Palaeohydrology of Laguna de Tagua Tagua
  (34° 30′ S) and moisture fluctuations in Central Chile for the last 46
  000 yr. J. Quat. Sci. 20, 625--641 (2005).
\item
  Herrera-Ossandón, M., Easton, G., Antinao, J. L. \& Forman, S. L. Late
  Quaternary glacier advances in the Andes of Santiago, central Chile,
  and paleoclimatic implications. Front. Earth Sci. 11, (2023).
\item
  Orellana, H., Latorre, C., García, J.-L. \& Lambert, F. Spatial
  analysis of paleoclimate variations based on proxy records in the
  south-central Andes (18°- 35° S) from 32 to 4 ka. Quat. Sci. Rev.~313,
  108174 (2023).
\item
  Garreaud, R. D. et al.~The 2010--2015 megadrought in central Chile:
  impacts on regional hydroclimate and vegetation. Hydrol. Earth Syst.
  Sci. 21, 6307--6327 (2017).
\item
  Garreaud, R. D., Vuille, M., Compagnucci, R. \& Marengo, J.
  Present-day South American climate. Palaeogeogr. Palaeoclimatol.
  Palaeoecol. 281, 180--195 (2009).
\item
  Gillett, Z. E., Hendon, H. H., Arblaster, J. M. \& Lim, E.-P. Tropical
  and Extratropical Influences on the Variability of the Southern
  Hemisphere Wintertime Subtropical Jet. J. Clim. 34, 4009--4022 (2021).
\item
  Lu, J., Chen, G. \& Frierson, D. M. W. Response of the Zonal Mean
  Atmospheric Circulation to El Niño versus Global Warming. J. Clim. 21,
  5835--5851 (2008).
\item
  Garreaud, R. Precipitation and Circulation Covariability in the
  Extratropics. J. Clim. 20, 4789--4797 (2007).
\item
  Campos, D. \& Rondanelli, R. ENSO-Related Precipitation Variability in
  Central Chile: The Role of Large Scale Moisture Transport. J. Geophys.
  Res. Atmospheres 128, e2023JD038671 (2023).
\item
  Collazo, S., García-Herrera, R. \& Barriopedro, D. Summer upper-level
  jets modulate the response of South American climate to ENSO. Clim.
  Dyn. 62, 1031--1054 (2024).
\item
  Moreno, P. I. \& León, A. L. Abrupt vegetation changes during the last
  glacial to Holocene transition in mid-latitude South America. J. Quat.
  Sci. 18, 787--800 (2003).
\item
  Sandweiss, D. H. et al.~Archaeological climate proxies and the
  complexities of reconstructing Holocene El Niño in coastal Peru. Proc.
  Natl. Acad. Sci. 117, 8271--8279 (2020).
\item
  Labarca, R. et al.~Taguatagua 3: A new late Pleistocene settlement in
  a highly suitable lacustrine habitat in central Chile (34°S). PLOS ONE
  19, e0302465 (2024).
\item
  Ugalde, P. C. et al.~The first peoples of the Atacama Desert lived
  among the trees: A 11,600- to 11,200-year-old grove and congregation
  site. Proc. Natl. Acad. Sci. 121, e2320506121 (2024).
\item
  Borrero, L. A. The Prehistoric Exploration and Colonization of
  Fuego-Patagonia. J. World Prehistory 13, 321--355 (1999).
\item
  Moy, C. M., Seltzer, G. O., Rodbell, D. T. \& Anderson, D. M.
  Variability of El Niño/Southern Oscillation activity at millennial
  timescales during the Holocene epoch. Nature 420, 162--165 (2002).
\item
  Moreno, P. I., Francois, J. P., Moy, C. M. \& Villa-Martínez, R.
  Covariability of the Southern Westerlies and atmospheric CO2 during
  the Holocene. Geology 38, 727--730 (2010).
\item
  Soteres, R. L. et al.~Glacier fluctuations in the northern Patagonian
  Andes (44°S) imply wind-modulated interhemispheric in-phase climate
  shifts during Termination 1. Sci. Rep.~12, 10842 (2022).
\item
  Villacís, L. A. et al.~A freshwater diatom perspective on the
  evolution of the southern westerlies for the past ∼14,000 years in
  southwestern Patagonia. Quat. Sci. Rev.~301, 107929 (2023).
\item
  Lamy, F. et al.~Holocene changes in the position and intensity of the
  southern westerly wind belt. Nat. Geosci. 3, 695--699 (2010).
\item
  Kilian, R. \& Lamy, F. A review of Glacial and Holocene paleoclimate
  records from southernmost Patagonia (49--55°S). Quat. Sci. Rev.~53,
  1--23 (2012).
\item
  Schneider, T., Bischoff, T. \& Haug, G. H. Migrations and dynamics of
  the intertropical convergence zone. Nature 513, 45--53 (2014).
\item
  Li, S. \& Liu, W. Deciphering the Migration of the Intertropical
  Convergence Zone During the Last Deglaciation. Geophys. Res. Lett. 49,
  e2022GL098806 (2022).
\item
  Kaiser, J. et al.~Orbital modulation of subtropical versus
  subantarctic moisture sources in the southeast Pacific mid-latitudes.
  Nat. Commun. 15, 7512 (2024).
\item
  Flores-Aqueveque, V., Rojas, M., Aguirre, C., Arias, P. A. \&
  González, C. South Pacific Subtropical High from the late Holocene to
  the end of the 21st century: insights from climate proxies and general
  circulation models. Clim. Past 16, 79--99 (2020).
\item
  Latorre, C., Betancourt, J. L. \& Arroyo, M. T. K. Late Quaternary
  vegetation and climate history of a perennial river canyon in the Río
  Salado basin (22°S) of Northern Chile. Quat. Res. 65, 450--466 (2006).
\item
  Pedro, J. B. et al.~The spatial extent and dynamics of the Antarctic
  Cold Reversal. Nat. Geosci. 9, 51--55 (2016).
\item
  Martin, L. C. P. et al.~Lake Tauca highstand (Heinrich Stadial 1a)
  driven by a southward shift of the Bolivian High. Sci. Adv. 4,
  eaar2514 (2018).
\item
  Prates, L. \& Perez, S. I. Late Pleistocene South American megafaunal
  extinctions associated with rise of Fishtail points and human
  population. Nat. Commun. 12, 2175 (2021).
\item
  Villavicencio, N. A. et al.~Combination of humans, climate, and
  vegetation change triggered Late Quaternary megafauna extinction in
  the Última Esperanza region, southern Patagonia, Chile. Ecography 39,
  125--140 (2016).
\item
  Timing and structure of the Younger Dryas event and its underlying
  climate dynamics \textbar{} PNAS.
  \url{https://www.pnas.org/doi/abs/10.1073/pnas.2007869117}.
\item
  Alley, R. B. The Younger Dryas cold interval as viewed from central
  Greenland. Quat. Sci. Rev.~19, 213--226 (2000).
\item
  Moreno-Mayar, J. V. et al.~Early human dispersals within the Americas.
  Science 362, (2018).
\item
  Liu, Z. et al.~Transient Simulation of Last Deglaciation with a New
  Mechanism for Bølling-Allerød Warming. Science 325, 310--314 (2009).
\item
  Heusser, C. J. Quaternary Pollen Record from Laguna de Tagua Tagua,
  Chile. Science 219, 1429--1432 (1983).
\item
  Varela, J. Geología del Cuaternario de Laguna de Taguatagua (Provincia
  de O'Higgins). in Congreso Geológico Chileno D81--D114 (1976).
\item
  Bird, D. et al.~p3k14c, a synthetic global database of archaeological
  radiocarbon dates. Sci. Data 9, 27 (2022).
\item
  Amann, B., Bertrand, S., Alvarez-Garreton, C. \& Reid, B. Seasonal
  Variations in Fjord Sediment Grain Size: A Pre-requisite for
  Hydrological and Climate Reconstructions in Partially Glacierized
  Watersheds (Baker River, Patagonia). J. Geophys. Res. Earth Surf. 127,
  e2021JF006391 (2022).
\item
  Sun, D. et al.~Grain-size distribution function of polymodal sediments
  in hydraulic and aeolian environments, and numerical partitioning of
  the sedimentary components. Sediment. Geol. 152, 263--277 (2002).
\item
  Vandenberghe, J., Sun, Y., Wang, X., Abels, H. A. \& Liu, X.
  Grain-size characterization of reworked fine-grained aeolian deposits.
  Earth-Sci. Rev.~177, 43--52 (2018).
\item
  Vandenberghe, J. Grain size of fine-grained windblown sediment: A
  powerful proxy for process identification. Earth-Sci. Rev.~121, 18--30
  (2013).
\item
  Dansgaard, W. Stable isotopes in precipitation. Tellus 16, 436--468
  (1964).
\item
  Rozanski, K., Araguás‐Araguás, L. \& Gonfiantini, R. Isotopic patterns
  in modern global precipitation. Clim. Change Cont. Isot. Rec. 78,
  1--36 (1993).
\item
  Gonfiantini, R., Roche, M.-A., Olivry, J.-C., Fontes, J.-C. \& Zuppi,
  G. M. The altitude effect on the isotopic composition of tropical
  rains. Chem. Geol. 181, 147--167 (2001).
\item
  Zielhofer, C. et al.~Millennial-scale fluctuations in Saharan dust
  supply across the decline of the African Humid Period. Quat. Sci.
  Rev.~171, 119--135 (2017).
\item
  Montané, J. Fechado del nivel superior de Tagua Tagua. Not. Mens. Mus.
  Nac. Hist. Nat. 14, 9--10 (1969).
\item
  Fathi-Moghadam, M., Arman, A., Emamgholizadeh, S. \& Alikhani, A.
  Settling Properties of Cohesive Sediments in Lakes and Reservoirs. J.
  Waterw. Port Coast. Ocean Eng. 137, 204--209 (2011).
\item
  Ciolkosz, E. J. et al.~Genesis of soils and landscapes in the Ridge
  and Valley province of central Pennsylvania. Geomorphology 3, 245--261
  (1990).
\item
  Beal, E. J., House, C. H. \& Orphan, V. J. Manganese- and
  Iron-Dependent Marine Methane Oxidation. Science 325, 184--187 (2009).
\item
  Wang, X. et al.~Anaerobic microbial manganese oxidation and reduction:
  A critical review. Sci. Total Environ. 822, 153513 (2022).
\item
  Dietze, E. et al.~Sediment transport processes across the Tibetan
  Plateau inferred from robust grain-size end members in lake sediments.
  Clim. Past 10, 91--106 (2014).
\item
  Zhang, X. et al.~Unmixing grain-size distributions in lake sediments:
  a new method of endmember modeling using hierarchical clustering.
  Quat. Res. 89, 365--373 (2018).
\item
  Orton, G. J. \& Reading, H. G. Variability of deltaic processes in
  terms of sediment supply, with particular emphasis on grain size.
  Sedimentology 40, 475--512 (1993).
\item
  van IJzendoorn, C. O., Hallin, C., Reniers, A. J. H. M. \& de Vries,
  S. Modeling Multi-Fraction Coastal Aeolian Sediment Transport With
  Horizontal and Vertical Grain-Size Variability. J. Geophys. Res. Earth
  Surf. 128, e2023JF007155 (2023).
\item
  van der Vegt, H. et al.~Grain size fractionation by process-driven
  sorting in sandy to muddy deltas. Depositional Rec. 6, 217--235
  (2020).
\item
  Heusser, C. J. Pollen and Spores of Chile. Modern Types of the
  Pteridophyta, Gymnospermae, and Angiospermae. (University of Arizona
  Press, 1971).
\item
  Mathiasen, P., Venegas-González, A., Fresia, P. \& Premoli, A. C. A
  relic of the past: current genetic patterns of the palaeoendemic tree
  Nothofagus macrocarpa were shaped by climatic oscillations in central
  Chile. Ann. Bot. 126, 891--904 (2020).
\item
  León-Lobos, P. \& Ellis, R. H. Seed survival in Chilean Nothofagus in
  response to desiccation and storage. Seed Sci. Res. 15, 113--123
  (2005).
\item
  Reyes-Díaz, M., Alberdi, M., Piper, F., Bravo, L. A. \& Corcuera, L.
  J. Low temperature responses of Nothofagus dombeyi and Nothofagus
  nitida, two evergreen species from south central Chile. Tree Physiol.
  25, 1389--1398 (2005).
\item
  Broecker, W. S. The Glacial World According to Wally. (Eldigio Press,
  1995).
\item
  Quade, J. et al.~Paleowetlands and regional climate change in the
  central Atacama Desert, northern Chile. Quat. Res. 69, 343--360
  (2008).
\item
  Gayo, E. M. et al.~Late Quaternary hydrological and ecological changes
  in the hyperarid core of the northern Atacama Desert
  (\textasciitilde{} 21°S). Earth-Sci. Rev.~113, 120--140 (2012).
\item
  González-Pinilla, F. J. et al.~High- and low-latitude forcings drive
  Atacama Desert rainfall variations over the past 16,000 years. Sci.
  Adv. 7, eabg1333 (2021).
\item
  Lamy, F., Hebbeln, D. \& Wefer, G. High-Resolution Marine Record of
  Climatic Change in Mid-latitude Chile during the Last 28,000 Years
  Based on Terrigenous Sediment Parameters. Quat. Res. 51, 83--93
  (1999).
\item
  Novello, V. F. et al.~A high-resolution history of the South American
  Monsoon from Last Glacial Maximum to the Holocene. Sci. Rep.~7, 44267
  (2017).
\item
  Kaiser, J., Schefu\ss, E., Lamy, F., Mohtadi, M. \& Hebbeln, D.
  Glacial to Holocene changes in sea surface temperature and coastal
  vegetation in north central Chile: high versus low latitude forcing.
  Quat. Sci. Rev.~27, 2064--2075 (2008).
\item
  Singer, B., Hildreth, W. \& Vincze, Y. 40Ar/39Ar evidence for early
  deglaciation of the central Chilean Andes. Geophys. Res. Lett. 27,
  1663--1666 (2000).
\item
  Fernández-Navarro, H. et al.~10Be chronology of the Last Glacial
  Maximum and Termination in the Andes of central Chile: The record of
  the Universidad Glacier (34° S). Quat. Sci. Rev.~344, 108968 (2024).
\item
  Kim, J.-H., Schneider, R. R., Hebbeln, D., Müller, P. J. \& Wefer, G.
  Last deglacial sea-surface temperature evolution in the Southeast
  Pacific compared to climate changes on the South American continent.
  Quat. Sci. Rev.~21, 2085--2097 (2002).
\item
  Kaiser, J., Lamy, F. \& Hebbeln, D. A 70-kyr sea surface temperature
  record off southern Chile (Ocean Drilling Program Site 1233).
  Paleoceanography 20, (2005).
\item
  Fritz, S. C. DIATOM METHODS \textbar{} Salinity and Climate
  Reconstructions from Continental Lakes. in Encyclopedia of Quaternary
  Science (Second Edition) (eds.~Elias, S. A. \& Mock, C. J.) 507--515
  (Elsevier, Amsterdam, 2013).
  \url{doi:10.1016/B978-0-444-53643-3.00221-1}.
\item
  Stenger-Kovács, C. et al.~Diatom community response to inland water
  salinization: a review. Hydrobiologia 850, 4627--4663 (2023).
\item
  Jara-Arancio, P. et al.~Genetic Diversity and Population Structure of
  Jubaea chilensis, an Endemic and Monotype Gender from Chile, Based on
  SNP Markers. Plants 11, 1959 (2022).
\item
  Trepel, J. et al.~Meta-analysis shows that wild large herbivores shape
  ecosystem properties and promote spatial heterogeneity. Nat. Ecol.
  Evol. 8, 705--716 (2024).
\item
  Robinson, M. et al.~`Moving South': Late Pleistocene Plant
  Exploitation and the Importance of Palm in the Colombian Amazon.
  Quaternary 4, 26 (2021).
\item
  McCrea, J. M. On the Isotopic Chemistry of Carbonates and a
  Paleotemperature Scale. J. Chem. Phys. 18, 849--857 (1950).
\item
  Horton, T. W., Defliese, W. F., Tripati, A. K. \& Oze, C. Evaporation
  induced 18O and 13C enrichment in lake systems: A global perspective
  on hydrologic balance effects. Quat. Sci. Rev.~131, 365--379 (2016).
\item
  McCormack, J., Nehrke, G., Jöns, N., Immenhauser, A. \& Kwiecien, O.
  Refining the interpretation of lacustrine carbonate isotope records:
  Implications of a mineralogy-specific Lake Van case study. Chem. Geol.
  513, 167--183 (2019).
\item
  Scott, W. P. et al.~Lake water based isoscape in central-south Chile
  reflects meteoric water. Sci. Rep.~11, 8725 (2021).
\item
  Terzer-Wassmuth, S., Wassenaar, L. I., Welker, J. M. \&
  Araguás-Araguás, L. J. Improved high-resolution global and
  regionalized isoscapes of δ18O, δ2H and d-excess in precipitation.
  Hydrol. Process. 35, e14254 (2021).
\item
  Benito, X., Trobajo, R. \& Ibáñez, C. Benthic diatoms in a
  Mediterranean delta: ecological indicators and a conductivity transfer
  function for paleoenvironmental studies. J. Paleolimnol. 54, 171--188
  (2015).
\item
  Starratt, S. W., Kirby, M. E. \& Glover, K. Diatom Record of Holocene
  Moisture Variability in the San Bernardino Mountains, California, USA.
  in Limnogeology: Progress, Challenges and Opportunities: A Tribute to
  Elizabeth Gierlowski-Kordesch (eds.~Rosen, M. R., Finkelstein, D. B.,
  Park Boush, L. \& Pla-Pueyo, S.) 329--365 (Springer International
  Publishing, Cham, 2021). \url{doi:10.1007/978-3-030-66576-0_11}.
\item
  Prates, L., Politis, G. G. \& Perez, S. I. Rapid radiation of humans
  in South America after the last glacial maximum: A radiocarbon-based
  study. PLOS ONE 15, e0236023 (2020).
\item
  Montané, J. Paleo-Indian Remains from Laguna de Tagua Tagua, Central
  Chile. Science 161, 1137--1138 (1968).
\item
  Méndez, C. \& Jackson, D. Terminal pleistocene lithic technology and
  use of space in Central Chile. Chungara Rev.~Antropol. Chil. 47,
  53--65 (2015).
\item
  Núñez, L., Varela, J., Casamiquela, R. \& Villagrán, C. Reconstruccion
  Multidisciplinaria de la Ocupacion Prehistorica de Quereo, Centro de
  Chile. Lat. Am. Antiq. 5, 99--118 (1994).
\item
  Lizama-Catalán, Á. \& Labarca, R. Who eats What: Unravelling a complex
  taphonomic scenario in the lacustrine deposits of the late Pleistocene
  archaeological site, Taguatagua 1, central Chile. Quat. Sci. Rev.~300,
  107831 (2023).
\item
  Rein, B. et al.~El Niño variability off Peru during the last 20,000
  years. Paleoceanography 20, (2005).
\item
  Jenny, B. et al.~Early to Mid-Holocene Aridity in Central Chile and
  the Southern Westerlies: The Laguna Aculeo Record (34°S). Quat. Res.
  58, 160--170 (2002).
\item
  Maldonado, A. et al.~Early Holocene climate change and human
  occupation along the semiarid coast of north-central Chile. J. Quat.
  Sci. 25, 985--988 (2010).
\item
  Ortega, C., Vargas, G., Rutllant, J. A., Jackson, D. \& Méndez, C.
  Major hydrological regime change along the semiarid western coast of
  South America during the early Holocene. Quat. Res. 78, 513--527
  (2012).
\item
  Frugone-Álvarez, M. et al.~Volcanism and climate change as drivers in
  Holocene depositional dynamic of Laguna del Maule (Andes of central
  Chile -- 36°S). Clim Past 16, 1097--1125 (2020).
\item
  Zhu, F. et al.~Climate models can correctly simulate the continuum of
  global-average temperature variability. Proc. Natl. Acad. Sci. 116,
  8728--8733 (2019).
\item
  Garreaud, R. D. et al.~The Central Chile Mega Drought (2010--2018): A
  climate dynamics perspective. Int. J. Climatol. 40, 421--439 (2020).
\item
  Jenny, B., Wilhelm, D. \& Valero-Garcés, B. The Southern Westerlies in
  Central Chile: Holocene precipitation estimates based on a water
  balance model for Laguna Aculeo (33 50′ S). Clim. Dyn. 20, 269--280
  (2003).
\item
  Tagle, B. \& del Río, C. Informe de prospecciones arqueológicas y
  excavación de túmulos funerarios en las provincias de Colchagua y
  Cachapoal. Inf. Al Cons. Monum. Nac. Ms (2005).
\item
  Barberena, R., Méndez, C. \& de Porras, M. E. Zooming out from
  archaeological discontinuities: The meaning of mid-Holocene temporal
  troughs in South American deserts. J. Anthropol. Archaeol. 46, 68--81
  (2017).
\item
  Núñez, L. \& Grosjean, M. Biodiversity and Human Impact During the
  Last 11,000 Years in North-Central Chile. in How Landscapes Change:
  Human Disturbance and Ecosystem Fragmentation in the Americas
  (eds.~Bradshaw, G. A. \& Marquet, P. A.) 7--17 (Springer, Berlin,
  Heidelberg, 2003). \url{doi:10.1007/978-3-662-05238-9_2}.
\item
  Magri, D. \& Di Rita, F. Archaeopalynological Preparation Techniques.
  in Plant Microtechniques and Protocols (eds.~Yeung, E. C. T.,
  Stasolla, C., Sumner, M. J. \& Huang, B. Q.) 495--506 (Springer
  International Publishing, Cham, 2015).
  \url{doi:10.1007/978-3-319-19944-3_27}.
\item
  Van Geel, B. Non-pollen palynomorphs. in Tracking environmental change
  using lake sediments: Terrestrial, algal, and siliceous indicators
  (eds.~Smol, J. P., Birks, H. J. B., Last, W. M., Bradley, R. S. \&
  Alverson, K.) 99--119 (Springer Netherlands, Dordrecht, 2001).
  \url{doi:10.1007/0-306-47668-1_6}.
\item
  Van Geel, B. et al.~Environmental reconstruction of a Roman Period
  settlement site in Uitgeest (The Netherlands), with special reference
  to coprophilous fungi. J. Archaeol. Sci. 30, 873--883 (2003).
\item
  Piperno, D. R. Phytolith Analysis. An Archaeological and Geological
  Perspective. 5 - laboratory techniques. in Phytolith analysis
  (ed.~Piperno, D. R.) 119--130 (Academic Press, San Diego, 1988).
  \url{doi:https://doi.org/10.1016/B978-0-12-557175-3.50009-X}.
\item
  Piperno, D. R. Phytoliths. In A Comprehensive Guide for Archaeologists
  and Palaeoecologists. (Altamira Press, USA, New York, 2006).
\item
  International Committee for Phytolith Taxonomy (ICPT). International
  Code for Phytolith Nomenclature (ICPN) 2.0. Ann. Bot. 124, 189--199
  (2019).
\item
  Morales, E. A., Siver, P. A. \& Trainor, F. R. Identification of
  diatoms (Bacillariophyceae) during ecological assessments: Comparison
  between Light Microscopy and Scanning Electron Microscopy techniques.
  Proc. Acad. Nat. Sci. Phila. 151, 95--103 (2001).
\item
  Hogg, A. G. et al.~SHCal20 Southern Hemisphere Calibration, 0--55,000
  Years cal BP. Radiocarbon 62, 759--778 (2020).
\item
  McKay, N. P., Emile-Geay, J. \& Khider, D. geoChronR -- an R package
  to model, analyze, and visualize age-uncertain data. Geochronology 3,
  149--169 (2021).
\item
  Blaauw, M. \& Christen, J. A. Flexible paleoclimate age-depth models
  using an autoregressive gamma process. Bayesian Anal. 6, 457--474
  (2011).
\item
  Parnell, A. C., Haslett, J., Allen, J. R. M., Buck, C. E. \& Huntley,
  B. A flexible approach to assessing synchroneity of past events using
  Bayesian reconstructions of sedimentation history. Quat. Sci. Rev.~27,
  1872--1885 (2008).
\item
  Dietze, E. \& Dietze, M. Grain-size distribution unmixing using the R
  package EMMAgeo. EG Quat. Sci. J. 68, 29--46 (2019).
\item
  Schnurrenberger, D., Russell, J. \& Kelts, K. Classification of
  lacustrine sediments based on sedimentary components. J. Paleolimnol.
  29, 141--154 (2003).
\item
  Hall, B. Facies classification using machine learning. Lead. Edge 35,
  906--909 (2016).
\item
  Pedregosa, F. et al.~Scikit-learn: Machine Learning in Python. J.
  Mach. Learn. Res. 12, 2825--2830 (2011).
\item
  Khider, D., Emile-Geay, J., Zhu, F., James, A., Landers, J., Ratnakar,
  V. \& Gil, Y. Pyleoclim: Paleoclimate timeseries analysis and
  visualization with Python. Paleoceanogr. Paleoclimatol. 37,
  e2022PA004509 (2022).
\end{enumerate}

\clearpage

\begin{table}[!hbtp]
\caption{Sedimentological, compositional characteristics and depositional environment of Tagua-Tagua lacustrine facies.}
\resizebox*{0.98\textwidth}{!}{
\begin{tabular}{p{2.5cm} p{5.5cm} p{5.5cm} p{5.2cm}}
\hline
\multicolumn{4}{l}{\textbf{1. Reworked diatomaceous mud soils}} \\
\hline \\
\textbf{Facies} & \textbf{Sedimentological properties} & \textbf{Composition and geochemistry} & \textbf{Depositional environment} \\ [0.5ex] 
\hline
\textbf{S1 (0-23 cm)} & Massive brown soil of fine silt (40 $\%$) to coarse clay (10 $\%$) with interbedded gravel clasts & High TOC (3.1 $\%$) and TN (0.5 $\%$) and low TIC (0.01 $\%$) and TS (0.05 $\%$). High $\delta^{13}{C}_{CO_{2}}$ and $\delta^{18}{O}_{CO_{2}}$ values ($\textrm{-21 \textperthousand}$ and $\textrm{-11 \textperthousand}$, respectively) & Soil with modern edaphic features soil with\\ [0.9ex]
\hline
\textbf{S2 (23-38 cm)} & Massive gray-brown to whitish-brown soil of coarse clay (25 $\%$) to coarse silt (40 $\%$). Presence of stratified (fissile) diatomaceous mud with charcoal fragments. & High TOC (2.3 $\%$) and TN (0.5 $\%$) and low TIC (0.04 $\%$) and TS (0.05 $\%$). High $\delta^{13}{C}_{CO_{2}}$ and $\delta^{18}{O}_{CO_{2}}$ values ($\textrm{-15 \textperthousand}$ and $\textrm{-7.8 \textperthousand}$, respectively) & Reworked fluvio-lacustrine sediment with modern edaphic features \\ [0.5ex]
\hline
\multicolumn{4}{l}{\textbf{2. Banded, diatomaceous, organic-poor silts}}\\  
\hline \\ 
\textbf{Facies} & \textbf{Sedimentological properties} & \textbf{Composition and geochemistry} & \textbf{Depositional environment} \\ [0.5ex] 
\hline
\textbf{L3 (38-79 cm)} & Coarse silt (30-40 $\%$) to fine sand (30 $\%$), with shell fragments of gastropods and bivalve arranged at random. Diatomaceous mud well consolidated fragments of between 0.1 to 0.5 cm and charcoal fragments little abundant & Lower TOC (1.2 $\%$), TN ($0.1$ $\%$) and TS ($0.1$ $\%$) and hight TIC ($> $$4$ $\%$). Lower $\delta^{13}{C}_{CO_{2}}$ and $\delta^{18}{O}_{CO_{2}}$ values ($\textrm{-1.15 \textperthousand}$ and $\textrm{-1.17 \textperthousand}$, respectively) & Littoral environments dominated by inorganic deposition, represent alternation from littoral to deeper environments\\ [0.5ex]
\hline
\textbf{D5} & Massive to banded, mm- to cm - thick layers of brown, macrophyte-rich coarse silt.& TS (1.4 $\%$) and presence of calcite (TIC = $0.1$$\%$). Higher TOC/TN. Moderate but fluctuating silicate-related elements (K, Ti, Rb, Sr and Zr), high Br/Ti, Ca/Ti, and S/Ti. High Mn values, and low Fe/Mn ratios.& \\ [0.5ex]
\hline
\multicolumn{4}{l}{\textbf{Laminated diatomaceous organic-rich silts(TOC 2-3.5$\%$; BioSi ca. 20 $\%$)}}\\ [0.7ex] 
\hline \hline
\textbf{Facies} & \textbf{Sedimentological properties} & \textbf{Composition and geochemistry} & \textbf{Depositional environment} \\ [0.7ex] 
\hline
\textbf{D6} &Laminated, dark brown, organic-rich silt composed of three main laminae type: a) greenish grey, diatomaceous medium to fine silt; b) dark brown, organic-rich fine silt with diffuse upper and lower boundaries and c) brown, macrophyte–rich, coarse silt, with sharp, erosional lower boundaries. Facies D6b contains more frequent endogenic calcite crystals.& High TOC (3.1$\%$), BioSi (19.7$\%$) and TS (1.4$\%$, peaks up to 6$\%$ in D6c) and common presence of calcite (TIC $>=$ 0.1). Moderate but fluctuating silicate-related elements (K, Ti, Rb, Sr and Zr); high Br/Ti, Ca/Ti, and S/Ti. High Mn values and low Fe/Mn ratios.& Distal, frequently oxic environments dominated by organic deposition and higher clastic input. Some periods with endogenic calcite formation\\ [0.5ex]
\hline
\end{tabular}}
\end{table}

\end{document}
